% =============================================================================
%  spike-gap.tex — เอกสารฉบับที่ 3 ของชุด "สัญญา tRPC" ของโครงการ ULMs
%
%  ระยะห่างระหว่างโค้ดจริงในรีโปกับ spike ที่ต่อสำเร็จแล้ว
%  เดินตามเส้นทางข้อมูล 5 สถานี — แต่ละสถานีบอก 3 อย่าง:
%      ทีมมีอะไรตอนนี้ · spike เปลี่ยนอะไร · ยังพังตรงไหน
%
%  ต้องคอมไพล์ด้วย XeLaTeX เท่านั้น (เอกสารเป็นภาษาไทย)
%      cd docs && latexmk spike-gap.tex        → build/spike-gap.pdf
%
%  ที่มาของข้อเท็จจริงทุกข้อ — ดูภาคผนวก ง ซึ่งอ้างไฟล์:บรรทัดครบทุกข้อ
%  สรุปแหล่งหลัก:
%    · frontend/src/**                          สภาพ frontend ของทีม ณ 10 ส.ค. 2569
%    · frontend/package.json                     zod ^3.23.8
%    · backend/src/**                            backend ของทีม (ยังไม่มี auth)
%    · backend/prisma/schema.prisma              33 model, 0 enum, ไม่มีตาราง Session
%    · spike-option-b/frontend-real/**           สำเนาที่ต่อ backend สำเร็จ
%    · spike-option-b/backend/src/**             backend ของ spike
%    · spike-option-b/browser-check/shots/*.png  ภาพหน้าจอจากการทดสอบจริง
%
%  รูปที่วาดด้วย TikZ = ภาพอธิบาย · รูปที่เป็น .png = ผลการทดสอบจริง
% =============================================================================

% =============================================================================
%  trpc-preamble.tex — preamble ร่วมของเอกสารชุด "สัญญา tRPC" ของโครงการ ULMs
%
%  ใช้โดย:  trpc-guide.tex    (ฉบับที่ 1 — หลักการ tRPC + ความเสี่ยงใน repo นี้)
%           trpc-meeting.tex  (ฉบับที่ 2 — วาระประชุมเรียงตามลำดับความสำคัญ)
%
%  ทำไมไม่ใช้ ulms-preamble.tex ที่มีอยู่แล้ว
%  ----------------------------------------
%  ulms-preamble.tex ตั้งใจให้เป็น "โทนเดียว (monochrome)" สำหรับเอกสารที่ต้อง
%  ส่งอาจารย์/พิมพ์ขาวดำ แต่เอกสารสอนสองฉบับนี้ใช้ "สีเชิงความหมาย" เป็นเครื่องมือ
%  สอนหลัก (สีเดียวกัน = แนวคิดเดียวกัน ทั้งในเนื้อความ ในไดอะแกรม และในตาราง)
%  จึงต้องมี preamble ของตัวเอง — ไม่ใช่ก๊อปมาแก้สีทับ
%
%  เอกสารที่เรียกใช้ต้องทำสองอย่างหลัง % =============================================================================
%  trpc-preamble.tex — preamble ร่วมของเอกสารชุด "สัญญา tRPC" ของโครงการ ULMs
%
%  ใช้โดย:  trpc-guide.tex    (ฉบับที่ 1 — หลักการ tRPC + ความเสี่ยงใน repo นี้)
%           trpc-meeting.tex  (ฉบับที่ 2 — วาระประชุมเรียงตามลำดับความสำคัญ)
%
%  ทำไมไม่ใช้ ulms-preamble.tex ที่มีอยู่แล้ว
%  ----------------------------------------
%  ulms-preamble.tex ตั้งใจให้เป็น "โทนเดียว (monochrome)" สำหรับเอกสารที่ต้อง
%  ส่งอาจารย์/พิมพ์ขาวดำ แต่เอกสารสอนสองฉบับนี้ใช้ "สีเชิงความหมาย" เป็นเครื่องมือ
%  สอนหลัก (สีเดียวกัน = แนวคิดเดียวกัน ทั้งในเนื้อความ ในไดอะแกรม และในตาราง)
%  จึงต้องมี preamble ของตัวเอง — ไม่ใช่ก๊อปมาแก้สีทับ
%
%  เอกสารที่เรียกใช้ต้องทำสองอย่างหลัง % =============================================================================
%  trpc-preamble.tex — preamble ร่วมของเอกสารชุด "สัญญา tRPC" ของโครงการ ULMs
%
%  ใช้โดย:  trpc-guide.tex    (ฉบับที่ 1 — หลักการ tRPC + ความเสี่ยงใน repo นี้)
%           trpc-meeting.tex  (ฉบับที่ 2 — วาระประชุมเรียงตามลำดับความสำคัญ)
%
%  ทำไมไม่ใช้ ulms-preamble.tex ที่มีอยู่แล้ว
%  ----------------------------------------
%  ulms-preamble.tex ตั้งใจให้เป็น "โทนเดียว (monochrome)" สำหรับเอกสารที่ต้อง
%  ส่งอาจารย์/พิมพ์ขาวดำ แต่เอกสารสอนสองฉบับนี้ใช้ "สีเชิงความหมาย" เป็นเครื่องมือ
%  สอนหลัก (สีเดียวกัน = แนวคิดเดียวกัน ทั้งในเนื้อความ ในไดอะแกรม และในตาราง)
%  จึงต้องมี preamble ของตัวเอง — ไม่ใช่ก๊อปมาแก้สีทับ
%
%  เอกสารที่เรียกใช้ต้องทำสองอย่างหลัง \input{trpc-preamble}:
%    1) \hypersetup{pdftitle={...}}
%    2) \begin{document}
%
%  Build:  cd docs && latexmk trpc-guide.tex     (.latexmkrc บังคับ xelatex ให้แล้ว)
%          PDF ออกที่ docs/build/ แล้วค่อยคัดลอกไป docs/report/
% =============================================================================

\documentclass[11pt,a4paper]{article}

% ---------------------------------------------------------------------
%  FONTS
%
%  ต้องใช้ XeLaTeX เท่านั้น — pdflatex เรียงพิมพ์ภาษาไทยไม่ได้เลย (ไม่ใช่ "ได้แต่ไม่สวย")
%
%  ใช้ตระกูล TLWG ที่มีทั้งอักษรไทยและละตินอยู่ในไฟล์เดียว จงใจเลี่ยงการจับคู่
%  ฟอนต์ละตินกับฟอนต์ไทยล้วนผ่าน ucharclasses เพราะ ucharclasses สลับฟอนต์ตาม
%  "บล็อกยูนิโคด" ของตัวอักษร โดยไม่รู้จักขอบเขตของ TeX group → ข้อความละตินที่
%  ตามหลัง \textbf{...} จะยังค้างอยู่ในฟอนต์ไทยซึ่งไม่มีสัญลักษณ์ละติน แล้ว
%  "หายไปจาก PDF เงียบ ๆ" โดยไม่มี warning
% ---------------------------------------------------------------------
\usepackage{fontspec}

\setmainfont{Laksaman}[Ligatures=TeX]          % ไทย+ละติน — เนื้อความ
\setsansfont{Garuda}[Ligatures=TeX]            % ไทย+ละติน หนักกว่า — หัวข้อ/ป้ายในไดอะแกรม
\setmonofont{DejaVu Sans Mono}[Scale=0.84]     % 0.84 ทำให้ x-height ใกล้เคียง TLWG

% --- การตัดบรรทัดภาษาไทย ---
% ภาษาไทยเขียนติดกันไม่มีช่องว่าง TeX จึงหาจุดตัดบรรทัดเองไม่ได้ และจะปล่อยให้
% บรรทัดล้นออกนอกหน้ากระดาษ สองบรรทัดนี้ยกงานให้ ICU ตัดคำแบบพจนานุกรม
% ค่า stretch ต้องมากกว่าศูนย์อย่างมีนัยสำคัญ ไม่งั้น TeX แทบไม่มีอิสระในการจัด
% บรรทัดไทยและจะออกมาเป็น overfull box แทนที่จะยืด/หดช่องไฟ
\XeTeXlinebreaklocale "th"
\XeTeXlinebreakskip = 0pt plus 0.15em minus 0.02em
\emergencystretch=2.2em

% หัวข้อ ป้ายในไดอะแกรม หัวตาราง header/footer ใช้ sans — เนื้อความไม่ใช้
% นี่คือสิ่งที่แยก "โครงสร้างสำหรับกวาดสายตา" ออกจาก "เนื้อหาสำหรับอ่าน"
\newcommand{\sansall}{\sffamily}

% ---------------------------------------------------------------------
%  LAYOUT
% ---------------------------------------------------------------------
\usepackage[a4paper,margin=2.3cm,top=2.5cm,bottom=2.5cm,headsep=12pt]{geometry}
\usepackage{amsmath,amssymb}
\usepackage{graphicx}
\graphicspath{{figures/}}
\usepackage{booktabs}
\usepackage{array}
\usepackage{longtable}
\usepackage{multirow}
\usepackage{xcolor}
\usepackage{tikz}
\usetikzlibrary{arrows.meta,positioning,calc,shapes.geometric,shapes.misc,
                fit,backgrounds,decorations.pathmorphing}
\usepackage[most]{tcolorbox}
\usepackage[font=small,labelfont=bf,width=0.95\textwidth]{caption}
\usepackage{fancyhdr}
\usepackage{enumitem}
\usepackage{titlesec}
\usepackage{float}        % figure[H] — ตรึงรูปไว้ข้างเนื้อความที่อธิบายมัน
\usepackage{microtype}
\usepackage{listings}
\usepackage[colorlinks=true,allcolors=ctr,breaklinks=true,
            bookmarksnumbered=true]{hyperref}

% ภาษาไทยต้องการระยะบรรทัดมากกว่าละติน สระบนซ้อนวรรณยุกต์ และมีสระล่างด้วย
% ที่ระยะปกติจะชนกันข้ามบรรทัด — 1.28 คือค่าที่ทดสอบแล้ว
\linespread{1.28}
\setlength{\parskip}{0.35em}
\setlength{\parindent}{0pt}

% ---------------------------------------------------------------------
%  สี — ระบบสีเชิงความหมาย (semantic colour)
%
%  ฐานสีคือชุด Okabe–Ito ซึ่งปลอดภัยสำหรับผู้ที่มองสีบางคู่ไม่แยก (colour-blind safe)
%
%  สีสี่สีหลักถูก "ตั้งชื่อกลาง ๆ" ไว้ที่นี่ แล้วให้แต่ละเอกสารประกาศความหมายของ
%  ตัวเองในหน้า "เอกสารนี้อ่านอย่างไร":
%
%    trpc-guide.tex   → 4 ชั้นของระบบ:   สัญญา / backend / frontend / เครือข่าย
%    trpc-meeting.tex → 4 ระดับเร่งด่วน: P0 / P1 / P2 / P3
%
%  สองเอกสารใช้คนละความหมายได้ เพราะแต่ละฉบับ "ประกาศระบบสีของตัวเองไว้หน้าแรก"
%  และคงเส้นคงวาภายในเล่ม — สิ่งที่ห้ามคือเปลี่ยนความหมายกลางเล่ม
% ---------------------------------------------------------------------
\definecolor{ctr}{HTML}{0072B2}      % ฟ้า   — สัญญา/schema/type   · หรือ P2
\definecolor{beo}{HTML}{D55E00}      % ส้ม   — ฝั่ง backend         · หรือ P1
\definecolor{feg}{HTML}{009E73}      % เขียว — ฝั่ง frontend        · หรือ P3
\definecolor{ops}{HTML}{7B5EA7}      % ม่วง  — เครือข่าย/รันไทม์/ops

\definecolor{ink}{HTML}{1A2733}      % เข้ม — หัวข้อรอง
\definecolor{muted}{HTML}{5B6B7A}    % ข้อความรอง เส้นคั่น ป้ายกำกับ
\definecolor{rule}{HTML}{C9D2DB}     % เส้นบาง ขอบไดอะแกรม
\definecolor{wash}{HTML}{F1F6FA}     % พื้นกล่องฟ้า/กลาง
\definecolor{warnwash}{HTML}{FDF3E7} % พื้นกล่องส้ม
\definecolor{goodwash}{HTML}{E9F4F0} % พื้นกล่องเขียว
\definecolor{badwash}{HTML}{FBEDED}  % พื้นกล่องแดง
\definecolor{bad}{HTML}{B3261E}      % แดง — ผิด/อันตราย · หรือ P0
\definecolor{codebg}{HTML}{F7F9FB}   % พื้นบล็อกโค้ด

% ---------------------------------------------------------------------
%  HEADINGS
%  \part ใช้เลขอารบิก เพื่อให้อ้างอิงข้ามเล่มอ่านว่า "ภาค 3" ไม่ใช่ "ภาค III"
% ---------------------------------------------------------------------
\renewcommand{\thepart}{\arabic{part}}
\titleformat{\part}[display]
  {\sansall\huge\bfseries\color{ctr}}
  {\normalsize\color{muted}ภาค \thepart}{0.4em}{}
\titlespacing*{\part}{0pt}{0pt}{1.5em}

\titleformat{\section}
  {\sansall\Large\bfseries\color{ctr!88!black}}{\thesection}{0.7em}{}
\titleformat{\subsection}
  {\sansall\large\bfseries\color{ink}}{\thesubsection}{0.6em}{}
\titleformat{\subsubsection}
  {\sansall\normalsize\bfseries\color{muted}}{\thesubsubsection}{0.55em}{}
\titlespacing*{\section}{0pt}{1.9em}{0.7em}
\titlespacing*{\subsection}{0pt}{1.3em}{0.45em}
\titlespacing*{\subsubsection}{0pt}{1.0em}{0.35em}

% หัวกระดาษ: ชื่อเอกสารคงที่ทางซ้าย ชื่อหัวข้อปัจจุบันทางขวา ผู้อ่านจะรู้ตลอดว่าอยู่ตรงไหน
% \docshort ถูกกำหนดในไฟล์เอกสารแต่ละฉบับ
\providecommand{\docshort}{}
\pagestyle{fancy}
\fancyhf{}
\renewcommand{\headrulewidth}{0.4pt}
\fancyhead[L]{\footnotesize\sansall\color{muted}\docshort}
\fancyhead[R]{\footnotesize\sansall\color{muted}\leftmark}
\fancyfoot[C]{\footnotesize\color{muted}\thepage}
\renewcommand{\sectionmark}[1]{\markboth{#1}{}}

% รายการแบบกระชับ — ค่าเริ่มต้นของ LaTeX ดูหลวมเมื่ออยู่กับ \parskip
\setlist[itemize]{itemsep=1pt,topsep=3pt,parsep=0pt,leftmargin=1.4em}
\setlist[enumerate]{itemsep=1pt,topsep=3pt,parsep=0pt,leftmargin=1.6em}
\setlist[description]{itemsep=2pt,topsep=3pt,parsep=0pt}

% ---------------------------------------------------------------------
%  กล่องข้อความ 4 แบบ ความหมายตายตัว ห้ามคิดแบบที่ห้ากลางเล่ม
%
%  colbacktitle ต้องตั้งให้เท่ากับ colback มิฉะนั้น tcolorbox จะใช้ colframe
%  (สีเข้ม) เป็นพื้นหลังของหัวกล่อง แล้วตัวอักษรหัวกล่องสีเข้มจะอ่านไม่ออก
% ---------------------------------------------------------------------
\tcbset{guidebox/.style={
  enhanced, breakable, boxrule=0pt, leftrule=2.5pt, arc=1pt,
  left=8pt, right=8pt, top=6pt, bottom=6pt,
  fonttitle=\sansall\bfseries\small, titlerule=0pt,
  toptitle=1pt, bottomtitle=1pt}}

% ฟ้า — ใจความที่ต้องจำ / ประโยคที่ร้อยหลายบทเข้าด้วยกัน
\newtcolorbox{keybox}[1][]{guidebox, colback=wash, colframe=ctr,
  colbacktitle=wash, coltitle=ctr!85!black, #1}
% ส้ม — กับดัก หรือข้อสมมติที่ถ้าลืมแล้วข้อสรุปจะผิด
\newtcolorbox{warnbox}[1][]{guidebox, colback=warnwash, colframe=beo,
  colbacktitle=warnwash, coltitle=beo!80!black, #1}
% เขียว — สิ่งที่ตรวจสอบกับของจริงในรีโปแล้ว
\newtcolorbox{goodbox}[1][]{guidebox, colback=goodwash, colframe=feg,
  colbacktitle=goodwash, coltitle=feg!65!black, #1}
% แดง — วิธีที่ผิดจริง ๆ อย่าทำ
\newtcolorbox{badbox}[1][]{guidebox, colback=badwash, colframe=bad,
  colbacktitle=badwash, coltitle=bad, #1}

% ---------------------------------------------------------------------
%  SEMANTIC MARKUP
% ---------------------------------------------------------------------
% ศัพท์อังกฤษที่แทรกในประโยคภาษาไทย ฟอนต์เนื้อความครอบคลุมละตินอยู่แล้วจึงไม่ต้อง
% สลับฟอนต์ — เก็บแมโครไว้เป็น "เครื่องหมายเชิงความหมาย" เผื่อวันหลังอยากเปลี่ยน
% การแสดงผลทั้งเล่มพร้อมกัน
\newcommand{\en}[1]{#1}
% ชื่อไฟล์ ฟังก์ชัน หรือค่าที่ปรากฏในโค้ดจริง
% ไม่บังคับขนาดตัวอักษร ปล่อยให้รับขนาดจากบริบท (สำคัญในตารางที่ใช้ \footnotesize
% มิฉะนั้นโค้ดจะกลายเป็นตัวใหญ่กว่าข้อความรอบข้างแล้วดันคอลัมน์ล้น)
% ฟอนต์ mono ถูกย่อไว้ 0.84 แล้วใน \setmonofont จึงพอดีกับเนื้อความอยู่แล้ว
\newcommand{\code}[1]{{\ttfamily #1}}

% ชื่อไฟล์ที่มีอักษรไทยอยู่ในชื่อ — ห้ามใช้ \code เพราะ DejaVu Sans Mono ไม่มี
% สัญลักษณ์ไทย ตัวอักษรจะหายไปจาก PDF โดยขึ้นแค่ warning "Missing character"
% ใช้ Garuda (sans ที่มีทั้งไทยและละติน) แทน
\newcommand{\fnth}[1]{{\sffamily\small #1}}

% ป้ายชื่อแนวคิด — สีเดียวกับที่ใช้ในไดอะแกรมและตาราง
\newcommand{\stg}[2]{\textcolor{#1}{\textbf{\en{#2}}}}

% ---------------------------------------------------------------------
%  ไวยากรณ์ของไดอะแกรม — นิยามครั้งเดียว ทุกรูปจะดูเป็นชุดเดียวกัน
%  หมายเหตุ: node ที่มีข้อความไทยต้องเผื่อ inner ysep มากกว่าไดอะแกรมภาษาละติน
%  เพราะสระบน/ล่างกินที่แนวตั้ง
% ---------------------------------------------------------------------
\tikzset{
  blk/.style={draw=rule, line width=0.6pt, rounded corners=2.5pt,
              fill=wash, minimum height=9mm, inner xsep=6pt, inner ysep=5pt,
              align=center, font=\sansall\footnotesize},
  stage/.style={blk, fill=ctr!8, draw=ctr!45},
  bestage/.style={blk, fill=beo!8, draw=beo!45},
  festage/.style={blk, fill=feg!8, draw=feg!45},
  opstage/.style={blk, fill=ops!8, draw=ops!45},
  ar/.style={-{Stealth[length=2.2mm,width=1.8mm]}, line width=0.7pt,
             draw=muted},
  lbl/.style={font=\sansall\scriptsize, text=muted, align=center},
  tag/.style={font=\sansall\scriptsize\bfseries, align=center},
}

% ---------------------------------------------------------------------
%  ตาราง
%  คอลัมน์แบบชิดซ้าย: ข้อความจัดชิดขอบสองข้าง (justified) ในคอลัมน์แคบจะเปิด
%  ช่องไฟกว้างมาก และแย่เป็นพิเศษกับภาษาไทย เพราะ ICU ตัดไทยเป็นก้อนที่ใหญ่กว่า
%  คำภาษาอังกฤษ → ใช้ L{} แทน p{} ในทุกคอลัมน์ที่มีข้อความยาว
% ---------------------------------------------------------------------
\newcolumntype{L}[1]{>{\raggedright\arraybackslash}p{#1}}
\captionsetup{skip=6pt}

% ---------------------------------------------------------------------
%  บล็อกโค้ด
%
%  ข้อจำกัดที่ต้องรู้: listings อ่านไฟล์แบบไบต์ ไม่เข้าใจ UTF-8 หลายไบต์
%  → "ห้ามใส่ภาษาไทยในบล็อกโค้ด" คอมเมนต์ในโค้ดทุกอันเป็นภาษาอังกฤษ
%    คำอธิบายภาษาไทยให้อยู่นอกบล็อกเสมอ
% ---------------------------------------------------------------------
\lstdefinelanguage{TypeScript}{
  keywords={await, async, break, case, catch, class, const, continue, default,
            delete, do, else, enum, export, extends, false, finally, for, from,
            function, if, implements, import, in, instanceof, interface, let,
            new, null, of, private, public, readonly, return, super, switch,
            this, throw, true, try, type, typeof, undefined, var, void, while,
            yield, satisfies, as},
  keywordstyle=\color{ctr}\bfseries,
  ndkeywords={string, number, boolean, Date, Promise, z, Record, Array},
  ndkeywordstyle=\color{ops},
  sensitive=true,
  comment=[l]{//},
  morecomment=[s]{/*}{*/},
  commentstyle=\color{muted}\itshape,
  stringstyle=\color{feg!70!black},
  morestring=[b]',
  morestring=[b]",
  morestring=[b]`,
}

\lstset{
  language=TypeScript,
  basicstyle=\ttfamily\scriptsize,
  backgroundcolor=\color{codebg},
  frame=leftline,
  framerule=2pt,
  rulecolor=\color{rule},
  framesep=7pt,
  xleftmargin=10pt,
  breaklines=true,
  breakatwhitespace=true,
  columns=fullflexible,
  keepspaces=true,
  showstringspaces=false,
  aboveskip=0.9em,
  belowskip=0.9em,
  literate={->}{{$\rightarrow$}}2 {=>}{{=>}}2,
}

% ชื่อ float ภาษาไทย — ถ้าไม่ตั้ง จะได้ "Figure 3" อยู่ใต้เนื้อความภาษาไทย
\renewcommand{\figurename}{รูปที่}
\renewcommand{\tablename}{ตารางที่}
\renewcommand{\refname}{เอกสารอ้างอิง}
\renewcommand{\contentsname}{สารบัญ}
% หมายเหตุ: ห้ามใส่ \color ใน \contentsname เพราะ \tableofcontents ส่งมันผ่าน
% \MakeUppercase ซึ่งจะเปลี่ยน ctr เป็น CTR แล้ว xcolor จะฟ้องว่าไม่รู้จักสีนั้น
% ถ้าจะจัดสไตล์หัวสารบัญ ให้ทำผ่าน \titleformat{\section} แทน
:
%    1) \hypersetup{pdftitle={...}}
%    2) \begin{document}
%
%  Build:  cd docs && latexmk trpc-guide.tex     (.latexmkrc บังคับ xelatex ให้แล้ว)
%          PDF ออกที่ docs/build/ แล้วค่อยคัดลอกไป docs/report/
% =============================================================================

\documentclass[11pt,a4paper]{article}

% ---------------------------------------------------------------------
%  FONTS
%
%  ต้องใช้ XeLaTeX เท่านั้น — pdflatex เรียงพิมพ์ภาษาไทยไม่ได้เลย (ไม่ใช่ "ได้แต่ไม่สวย")
%
%  ใช้ตระกูล TLWG ที่มีทั้งอักษรไทยและละตินอยู่ในไฟล์เดียว จงใจเลี่ยงการจับคู่
%  ฟอนต์ละตินกับฟอนต์ไทยล้วนผ่าน ucharclasses เพราะ ucharclasses สลับฟอนต์ตาม
%  "บล็อกยูนิโคด" ของตัวอักษร โดยไม่รู้จักขอบเขตของ TeX group → ข้อความละตินที่
%  ตามหลัง \textbf{...} จะยังค้างอยู่ในฟอนต์ไทยซึ่งไม่มีสัญลักษณ์ละติน แล้ว
%  "หายไปจาก PDF เงียบ ๆ" โดยไม่มี warning
% ---------------------------------------------------------------------
\usepackage{fontspec}

\setmainfont{Laksaman}[Ligatures=TeX]          % ไทย+ละติน — เนื้อความ
\setsansfont{Garuda}[Ligatures=TeX]            % ไทย+ละติน หนักกว่า — หัวข้อ/ป้ายในไดอะแกรม
\setmonofont{DejaVu Sans Mono}[Scale=0.84]     % 0.84 ทำให้ x-height ใกล้เคียง TLWG

% --- การตัดบรรทัดภาษาไทย ---
% ภาษาไทยเขียนติดกันไม่มีช่องว่าง TeX จึงหาจุดตัดบรรทัดเองไม่ได้ และจะปล่อยให้
% บรรทัดล้นออกนอกหน้ากระดาษ สองบรรทัดนี้ยกงานให้ ICU ตัดคำแบบพจนานุกรม
% ค่า stretch ต้องมากกว่าศูนย์อย่างมีนัยสำคัญ ไม่งั้น TeX แทบไม่มีอิสระในการจัด
% บรรทัดไทยและจะออกมาเป็น overfull box แทนที่จะยืด/หดช่องไฟ
\XeTeXlinebreaklocale "th"
\XeTeXlinebreakskip = 0pt plus 0.15em minus 0.02em
\emergencystretch=2.2em

% หัวข้อ ป้ายในไดอะแกรม หัวตาราง header/footer ใช้ sans — เนื้อความไม่ใช้
% นี่คือสิ่งที่แยก "โครงสร้างสำหรับกวาดสายตา" ออกจาก "เนื้อหาสำหรับอ่าน"
\newcommand{\sansall}{\sffamily}

% ---------------------------------------------------------------------
%  LAYOUT
% ---------------------------------------------------------------------
\usepackage[a4paper,margin=2.3cm,top=2.5cm,bottom=2.5cm,headsep=12pt]{geometry}
\usepackage{amsmath,amssymb}
\usepackage{graphicx}
\graphicspath{{figures/}}
\usepackage{booktabs}
\usepackage{array}
\usepackage{longtable}
\usepackage{multirow}
\usepackage{xcolor}
\usepackage{tikz}
\usetikzlibrary{arrows.meta,positioning,calc,shapes.geometric,shapes.misc,
                fit,backgrounds,decorations.pathmorphing}
\usepackage[most]{tcolorbox}
\usepackage[font=small,labelfont=bf,width=0.95\textwidth]{caption}
\usepackage{fancyhdr}
\usepackage{enumitem}
\usepackage{titlesec}
\usepackage{float}        % figure[H] — ตรึงรูปไว้ข้างเนื้อความที่อธิบายมัน
\usepackage{microtype}
\usepackage{listings}
\usepackage[colorlinks=true,allcolors=ctr,breaklinks=true,
            bookmarksnumbered=true]{hyperref}

% ภาษาไทยต้องการระยะบรรทัดมากกว่าละติน สระบนซ้อนวรรณยุกต์ และมีสระล่างด้วย
% ที่ระยะปกติจะชนกันข้ามบรรทัด — 1.28 คือค่าที่ทดสอบแล้ว
\linespread{1.28}
\setlength{\parskip}{0.35em}
\setlength{\parindent}{0pt}

% ---------------------------------------------------------------------
%  สี — ระบบสีเชิงความหมาย (semantic colour)
%
%  ฐานสีคือชุด Okabe–Ito ซึ่งปลอดภัยสำหรับผู้ที่มองสีบางคู่ไม่แยก (colour-blind safe)
%
%  สีสี่สีหลักถูก "ตั้งชื่อกลาง ๆ" ไว้ที่นี่ แล้วให้แต่ละเอกสารประกาศความหมายของ
%  ตัวเองในหน้า "เอกสารนี้อ่านอย่างไร":
%
%    trpc-guide.tex   → 4 ชั้นของระบบ:   สัญญา / backend / frontend / เครือข่าย
%    trpc-meeting.tex → 4 ระดับเร่งด่วน: P0 / P1 / P2 / P3
%
%  สองเอกสารใช้คนละความหมายได้ เพราะแต่ละฉบับ "ประกาศระบบสีของตัวเองไว้หน้าแรก"
%  และคงเส้นคงวาภายในเล่ม — สิ่งที่ห้ามคือเปลี่ยนความหมายกลางเล่ม
% ---------------------------------------------------------------------
\definecolor{ctr}{HTML}{0072B2}      % ฟ้า   — สัญญา/schema/type   · หรือ P2
\definecolor{beo}{HTML}{D55E00}      % ส้ม   — ฝั่ง backend         · หรือ P1
\definecolor{feg}{HTML}{009E73}      % เขียว — ฝั่ง frontend        · หรือ P3
\definecolor{ops}{HTML}{7B5EA7}      % ม่วง  — เครือข่าย/รันไทม์/ops

\definecolor{ink}{HTML}{1A2733}      % เข้ม — หัวข้อรอง
\definecolor{muted}{HTML}{5B6B7A}    % ข้อความรอง เส้นคั่น ป้ายกำกับ
\definecolor{rule}{HTML}{C9D2DB}     % เส้นบาง ขอบไดอะแกรม
\definecolor{wash}{HTML}{F1F6FA}     % พื้นกล่องฟ้า/กลาง
\definecolor{warnwash}{HTML}{FDF3E7} % พื้นกล่องส้ม
\definecolor{goodwash}{HTML}{E9F4F0} % พื้นกล่องเขียว
\definecolor{badwash}{HTML}{FBEDED}  % พื้นกล่องแดง
\definecolor{bad}{HTML}{B3261E}      % แดง — ผิด/อันตราย · หรือ P0
\definecolor{codebg}{HTML}{F7F9FB}   % พื้นบล็อกโค้ด

% ---------------------------------------------------------------------
%  HEADINGS
%  \part ใช้เลขอารบิก เพื่อให้อ้างอิงข้ามเล่มอ่านว่า "ภาค 3" ไม่ใช่ "ภาค III"
% ---------------------------------------------------------------------
\renewcommand{\thepart}{\arabic{part}}
\titleformat{\part}[display]
  {\sansall\huge\bfseries\color{ctr}}
  {\normalsize\color{muted}ภาค \thepart}{0.4em}{}
\titlespacing*{\part}{0pt}{0pt}{1.5em}

\titleformat{\section}
  {\sansall\Large\bfseries\color{ctr!88!black}}{\thesection}{0.7em}{}
\titleformat{\subsection}
  {\sansall\large\bfseries\color{ink}}{\thesubsection}{0.6em}{}
\titleformat{\subsubsection}
  {\sansall\normalsize\bfseries\color{muted}}{\thesubsubsection}{0.55em}{}
\titlespacing*{\section}{0pt}{1.9em}{0.7em}
\titlespacing*{\subsection}{0pt}{1.3em}{0.45em}
\titlespacing*{\subsubsection}{0pt}{1.0em}{0.35em}

% หัวกระดาษ: ชื่อเอกสารคงที่ทางซ้าย ชื่อหัวข้อปัจจุบันทางขวา ผู้อ่านจะรู้ตลอดว่าอยู่ตรงไหน
% \docshort ถูกกำหนดในไฟล์เอกสารแต่ละฉบับ
\providecommand{\docshort}{}
\pagestyle{fancy}
\fancyhf{}
\renewcommand{\headrulewidth}{0.4pt}
\fancyhead[L]{\footnotesize\sansall\color{muted}\docshort}
\fancyhead[R]{\footnotesize\sansall\color{muted}\leftmark}
\fancyfoot[C]{\footnotesize\color{muted}\thepage}
\renewcommand{\sectionmark}[1]{\markboth{#1}{}}

% รายการแบบกระชับ — ค่าเริ่มต้นของ LaTeX ดูหลวมเมื่ออยู่กับ \parskip
\setlist[itemize]{itemsep=1pt,topsep=3pt,parsep=0pt,leftmargin=1.4em}
\setlist[enumerate]{itemsep=1pt,topsep=3pt,parsep=0pt,leftmargin=1.6em}
\setlist[description]{itemsep=2pt,topsep=3pt,parsep=0pt}

% ---------------------------------------------------------------------
%  กล่องข้อความ 4 แบบ ความหมายตายตัว ห้ามคิดแบบที่ห้ากลางเล่ม
%
%  colbacktitle ต้องตั้งให้เท่ากับ colback มิฉะนั้น tcolorbox จะใช้ colframe
%  (สีเข้ม) เป็นพื้นหลังของหัวกล่อง แล้วตัวอักษรหัวกล่องสีเข้มจะอ่านไม่ออก
% ---------------------------------------------------------------------
\tcbset{guidebox/.style={
  enhanced, breakable, boxrule=0pt, leftrule=2.5pt, arc=1pt,
  left=8pt, right=8pt, top=6pt, bottom=6pt,
  fonttitle=\sansall\bfseries\small, titlerule=0pt,
  toptitle=1pt, bottomtitle=1pt}}

% ฟ้า — ใจความที่ต้องจำ / ประโยคที่ร้อยหลายบทเข้าด้วยกัน
\newtcolorbox{keybox}[1][]{guidebox, colback=wash, colframe=ctr,
  colbacktitle=wash, coltitle=ctr!85!black, #1}
% ส้ม — กับดัก หรือข้อสมมติที่ถ้าลืมแล้วข้อสรุปจะผิด
\newtcolorbox{warnbox}[1][]{guidebox, colback=warnwash, colframe=beo,
  colbacktitle=warnwash, coltitle=beo!80!black, #1}
% เขียว — สิ่งที่ตรวจสอบกับของจริงในรีโปแล้ว
\newtcolorbox{goodbox}[1][]{guidebox, colback=goodwash, colframe=feg,
  colbacktitle=goodwash, coltitle=feg!65!black, #1}
% แดง — วิธีที่ผิดจริง ๆ อย่าทำ
\newtcolorbox{badbox}[1][]{guidebox, colback=badwash, colframe=bad,
  colbacktitle=badwash, coltitle=bad, #1}

% ---------------------------------------------------------------------
%  SEMANTIC MARKUP
% ---------------------------------------------------------------------
% ศัพท์อังกฤษที่แทรกในประโยคภาษาไทย ฟอนต์เนื้อความครอบคลุมละตินอยู่แล้วจึงไม่ต้อง
% สลับฟอนต์ — เก็บแมโครไว้เป็น "เครื่องหมายเชิงความหมาย" เผื่อวันหลังอยากเปลี่ยน
% การแสดงผลทั้งเล่มพร้อมกัน
\newcommand{\en}[1]{#1}
% ชื่อไฟล์ ฟังก์ชัน หรือค่าที่ปรากฏในโค้ดจริง
% ไม่บังคับขนาดตัวอักษร ปล่อยให้รับขนาดจากบริบท (สำคัญในตารางที่ใช้ \footnotesize
% มิฉะนั้นโค้ดจะกลายเป็นตัวใหญ่กว่าข้อความรอบข้างแล้วดันคอลัมน์ล้น)
% ฟอนต์ mono ถูกย่อไว้ 0.84 แล้วใน \setmonofont จึงพอดีกับเนื้อความอยู่แล้ว
\newcommand{\code}[1]{{\ttfamily #1}}

% ชื่อไฟล์ที่มีอักษรไทยอยู่ในชื่อ — ห้ามใช้ \code เพราะ DejaVu Sans Mono ไม่มี
% สัญลักษณ์ไทย ตัวอักษรจะหายไปจาก PDF โดยขึ้นแค่ warning "Missing character"
% ใช้ Garuda (sans ที่มีทั้งไทยและละติน) แทน
\newcommand{\fnth}[1]{{\sffamily\small #1}}

% ป้ายชื่อแนวคิด — สีเดียวกับที่ใช้ในไดอะแกรมและตาราง
\newcommand{\stg}[2]{\textcolor{#1}{\textbf{\en{#2}}}}

% ---------------------------------------------------------------------
%  ไวยากรณ์ของไดอะแกรม — นิยามครั้งเดียว ทุกรูปจะดูเป็นชุดเดียวกัน
%  หมายเหตุ: node ที่มีข้อความไทยต้องเผื่อ inner ysep มากกว่าไดอะแกรมภาษาละติน
%  เพราะสระบน/ล่างกินที่แนวตั้ง
% ---------------------------------------------------------------------
\tikzset{
  blk/.style={draw=rule, line width=0.6pt, rounded corners=2.5pt,
              fill=wash, minimum height=9mm, inner xsep=6pt, inner ysep=5pt,
              align=center, font=\sansall\footnotesize},
  stage/.style={blk, fill=ctr!8, draw=ctr!45},
  bestage/.style={blk, fill=beo!8, draw=beo!45},
  festage/.style={blk, fill=feg!8, draw=feg!45},
  opstage/.style={blk, fill=ops!8, draw=ops!45},
  ar/.style={-{Stealth[length=2.2mm,width=1.8mm]}, line width=0.7pt,
             draw=muted},
  lbl/.style={font=\sansall\scriptsize, text=muted, align=center},
  tag/.style={font=\sansall\scriptsize\bfseries, align=center},
}

% ---------------------------------------------------------------------
%  ตาราง
%  คอลัมน์แบบชิดซ้าย: ข้อความจัดชิดขอบสองข้าง (justified) ในคอลัมน์แคบจะเปิด
%  ช่องไฟกว้างมาก และแย่เป็นพิเศษกับภาษาไทย เพราะ ICU ตัดไทยเป็นก้อนที่ใหญ่กว่า
%  คำภาษาอังกฤษ → ใช้ L{} แทน p{} ในทุกคอลัมน์ที่มีข้อความยาว
% ---------------------------------------------------------------------
\newcolumntype{L}[1]{>{\raggedright\arraybackslash}p{#1}}
\captionsetup{skip=6pt}

% ---------------------------------------------------------------------
%  บล็อกโค้ด
%
%  ข้อจำกัดที่ต้องรู้: listings อ่านไฟล์แบบไบต์ ไม่เข้าใจ UTF-8 หลายไบต์
%  → "ห้ามใส่ภาษาไทยในบล็อกโค้ด" คอมเมนต์ในโค้ดทุกอันเป็นภาษาอังกฤษ
%    คำอธิบายภาษาไทยให้อยู่นอกบล็อกเสมอ
% ---------------------------------------------------------------------
\lstdefinelanguage{TypeScript}{
  keywords={await, async, break, case, catch, class, const, continue, default,
            delete, do, else, enum, export, extends, false, finally, for, from,
            function, if, implements, import, in, instanceof, interface, let,
            new, null, of, private, public, readonly, return, super, switch,
            this, throw, true, try, type, typeof, undefined, var, void, while,
            yield, satisfies, as},
  keywordstyle=\color{ctr}\bfseries,
  ndkeywords={string, number, boolean, Date, Promise, z, Record, Array},
  ndkeywordstyle=\color{ops},
  sensitive=true,
  comment=[l]{//},
  morecomment=[s]{/*}{*/},
  commentstyle=\color{muted}\itshape,
  stringstyle=\color{feg!70!black},
  morestring=[b]',
  morestring=[b]",
  morestring=[b]`,
}

\lstset{
  language=TypeScript,
  basicstyle=\ttfamily\scriptsize,
  backgroundcolor=\color{codebg},
  frame=leftline,
  framerule=2pt,
  rulecolor=\color{rule},
  framesep=7pt,
  xleftmargin=10pt,
  breaklines=true,
  breakatwhitespace=true,
  columns=fullflexible,
  keepspaces=true,
  showstringspaces=false,
  aboveskip=0.9em,
  belowskip=0.9em,
  literate={->}{{$\rightarrow$}}2 {=>}{{=>}}2,
}

% ชื่อ float ภาษาไทย — ถ้าไม่ตั้ง จะได้ "Figure 3" อยู่ใต้เนื้อความภาษาไทย
\renewcommand{\figurename}{รูปที่}
\renewcommand{\tablename}{ตารางที่}
\renewcommand{\refname}{เอกสารอ้างอิง}
\renewcommand{\contentsname}{สารบัญ}
% หมายเหตุ: ห้ามใส่ \color ใน \contentsname เพราะ \tableofcontents ส่งมันผ่าน
% \MakeUppercase ซึ่งจะเปลี่ยน ctr เป็น CTR แล้ว xcolor จะฟ้องว่าไม่รู้จักสีนั้น
% ถ้าจะจัดสไตล์หัวสารบัญ ให้ทำผ่าน \titleformat{\section} แทน
:
%    1) \hypersetup{pdftitle={...}}
%    2) \begin{document}
%
%  Build:  cd docs && latexmk trpc-guide.tex     (.latexmkrc บังคับ xelatex ให้แล้ว)
%          PDF ออกที่ docs/build/ แล้วค่อยคัดลอกไป docs/report/
% =============================================================================

\documentclass[11pt,a4paper]{article}

% ---------------------------------------------------------------------
%  FONTS
%
%  ต้องใช้ XeLaTeX เท่านั้น — pdflatex เรียงพิมพ์ภาษาไทยไม่ได้เลย (ไม่ใช่ "ได้แต่ไม่สวย")
%
%  ใช้ตระกูล TLWG ที่มีทั้งอักษรไทยและละตินอยู่ในไฟล์เดียว จงใจเลี่ยงการจับคู่
%  ฟอนต์ละตินกับฟอนต์ไทยล้วนผ่าน ucharclasses เพราะ ucharclasses สลับฟอนต์ตาม
%  "บล็อกยูนิโคด" ของตัวอักษร โดยไม่รู้จักขอบเขตของ TeX group → ข้อความละตินที่
%  ตามหลัง \textbf{...} จะยังค้างอยู่ในฟอนต์ไทยซึ่งไม่มีสัญลักษณ์ละติน แล้ว
%  "หายไปจาก PDF เงียบ ๆ" โดยไม่มี warning
% ---------------------------------------------------------------------
\usepackage{fontspec}

\setmainfont{Laksaman}[Ligatures=TeX]          % ไทย+ละติน — เนื้อความ
\setsansfont{Garuda}[Ligatures=TeX]            % ไทย+ละติน หนักกว่า — หัวข้อ/ป้ายในไดอะแกรม
\setmonofont{DejaVu Sans Mono}[Scale=0.84]     % 0.84 ทำให้ x-height ใกล้เคียง TLWG

% --- การตัดบรรทัดภาษาไทย ---
% ภาษาไทยเขียนติดกันไม่มีช่องว่าง TeX จึงหาจุดตัดบรรทัดเองไม่ได้ และจะปล่อยให้
% บรรทัดล้นออกนอกหน้ากระดาษ สองบรรทัดนี้ยกงานให้ ICU ตัดคำแบบพจนานุกรม
% ค่า stretch ต้องมากกว่าศูนย์อย่างมีนัยสำคัญ ไม่งั้น TeX แทบไม่มีอิสระในการจัด
% บรรทัดไทยและจะออกมาเป็น overfull box แทนที่จะยืด/หดช่องไฟ
\XeTeXlinebreaklocale "th"
\XeTeXlinebreakskip = 0pt plus 0.15em minus 0.02em
\emergencystretch=2.2em

% หัวข้อ ป้ายในไดอะแกรม หัวตาราง header/footer ใช้ sans — เนื้อความไม่ใช้
% นี่คือสิ่งที่แยก "โครงสร้างสำหรับกวาดสายตา" ออกจาก "เนื้อหาสำหรับอ่าน"
\newcommand{\sansall}{\sffamily}

% ---------------------------------------------------------------------
%  LAYOUT
% ---------------------------------------------------------------------
\usepackage[a4paper,margin=2.3cm,top=2.5cm,bottom=2.5cm,headsep=12pt]{geometry}
\usepackage{amsmath,amssymb}
\usepackage{graphicx}
\graphicspath{{figures/}}
\usepackage{booktabs}
\usepackage{array}
\usepackage{longtable}
\usepackage{multirow}
\usepackage{xcolor}
\usepackage{tikz}
\usetikzlibrary{arrows.meta,positioning,calc,shapes.geometric,shapes.misc,
                fit,backgrounds,decorations.pathmorphing}
\usepackage[most]{tcolorbox}
\usepackage[font=small,labelfont=bf,width=0.95\textwidth]{caption}
\usepackage{fancyhdr}
\usepackage{enumitem}
\usepackage{titlesec}
\usepackage{float}        % figure[H] — ตรึงรูปไว้ข้างเนื้อความที่อธิบายมัน
\usepackage{microtype}
\usepackage{listings}
\usepackage[colorlinks=true,allcolors=ctr,breaklinks=true,
            bookmarksnumbered=true]{hyperref}

% ภาษาไทยต้องการระยะบรรทัดมากกว่าละติน สระบนซ้อนวรรณยุกต์ และมีสระล่างด้วย
% ที่ระยะปกติจะชนกันข้ามบรรทัด — 1.28 คือค่าที่ทดสอบแล้ว
\linespread{1.28}
\setlength{\parskip}{0.35em}
\setlength{\parindent}{0pt}

% ---------------------------------------------------------------------
%  สี — ระบบสีเชิงความหมาย (semantic colour)
%
%  ฐานสีคือชุด Okabe–Ito ซึ่งปลอดภัยสำหรับผู้ที่มองสีบางคู่ไม่แยก (colour-blind safe)
%
%  สีสี่สีหลักถูก "ตั้งชื่อกลาง ๆ" ไว้ที่นี่ แล้วให้แต่ละเอกสารประกาศความหมายของ
%  ตัวเองในหน้า "เอกสารนี้อ่านอย่างไร":
%
%    trpc-guide.tex   → 4 ชั้นของระบบ:   สัญญา / backend / frontend / เครือข่าย
%    trpc-meeting.tex → 4 ระดับเร่งด่วน: P0 / P1 / P2 / P3
%
%  สองเอกสารใช้คนละความหมายได้ เพราะแต่ละฉบับ "ประกาศระบบสีของตัวเองไว้หน้าแรก"
%  และคงเส้นคงวาภายในเล่ม — สิ่งที่ห้ามคือเปลี่ยนความหมายกลางเล่ม
% ---------------------------------------------------------------------
\definecolor{ctr}{HTML}{0072B2}      % ฟ้า   — สัญญา/schema/type   · หรือ P2
\definecolor{beo}{HTML}{D55E00}      % ส้ม   — ฝั่ง backend         · หรือ P1
\definecolor{feg}{HTML}{009E73}      % เขียว — ฝั่ง frontend        · หรือ P3
\definecolor{ops}{HTML}{7B5EA7}      % ม่วง  — เครือข่าย/รันไทม์/ops

\definecolor{ink}{HTML}{1A2733}      % เข้ม — หัวข้อรอง
\definecolor{muted}{HTML}{5B6B7A}    % ข้อความรอง เส้นคั่น ป้ายกำกับ
\definecolor{rule}{HTML}{C9D2DB}     % เส้นบาง ขอบไดอะแกรม
\definecolor{wash}{HTML}{F1F6FA}     % พื้นกล่องฟ้า/กลาง
\definecolor{warnwash}{HTML}{FDF3E7} % พื้นกล่องส้ม
\definecolor{goodwash}{HTML}{E9F4F0} % พื้นกล่องเขียว
\definecolor{badwash}{HTML}{FBEDED}  % พื้นกล่องแดง
\definecolor{bad}{HTML}{B3261E}      % แดง — ผิด/อันตราย · หรือ P0
\definecolor{codebg}{HTML}{F7F9FB}   % พื้นบล็อกโค้ด

% ---------------------------------------------------------------------
%  HEADINGS
%  \part ใช้เลขอารบิก เพื่อให้อ้างอิงข้ามเล่มอ่านว่า "ภาค 3" ไม่ใช่ "ภาค III"
% ---------------------------------------------------------------------
\renewcommand{\thepart}{\arabic{part}}
\titleformat{\part}[display]
  {\sansall\huge\bfseries\color{ctr}}
  {\normalsize\color{muted}ภาค \thepart}{0.4em}{}
\titlespacing*{\part}{0pt}{0pt}{1.5em}

\titleformat{\section}
  {\sansall\Large\bfseries\color{ctr!88!black}}{\thesection}{0.7em}{}
\titleformat{\subsection}
  {\sansall\large\bfseries\color{ink}}{\thesubsection}{0.6em}{}
\titleformat{\subsubsection}
  {\sansall\normalsize\bfseries\color{muted}}{\thesubsubsection}{0.55em}{}
\titlespacing*{\section}{0pt}{1.9em}{0.7em}
\titlespacing*{\subsection}{0pt}{1.3em}{0.45em}
\titlespacing*{\subsubsection}{0pt}{1.0em}{0.35em}

% หัวกระดาษ: ชื่อเอกสารคงที่ทางซ้าย ชื่อหัวข้อปัจจุบันทางขวา ผู้อ่านจะรู้ตลอดว่าอยู่ตรงไหน
% \docshort ถูกกำหนดในไฟล์เอกสารแต่ละฉบับ
\providecommand{\docshort}{}
\pagestyle{fancy}
\fancyhf{}
\renewcommand{\headrulewidth}{0.4pt}
\fancyhead[L]{\footnotesize\sansall\color{muted}\docshort}
\fancyhead[R]{\footnotesize\sansall\color{muted}\leftmark}
\fancyfoot[C]{\footnotesize\color{muted}\thepage}
\renewcommand{\sectionmark}[1]{\markboth{#1}{}}

% รายการแบบกระชับ — ค่าเริ่มต้นของ LaTeX ดูหลวมเมื่ออยู่กับ \parskip
\setlist[itemize]{itemsep=1pt,topsep=3pt,parsep=0pt,leftmargin=1.4em}
\setlist[enumerate]{itemsep=1pt,topsep=3pt,parsep=0pt,leftmargin=1.6em}
\setlist[description]{itemsep=2pt,topsep=3pt,parsep=0pt}

% ---------------------------------------------------------------------
%  กล่องข้อความ 4 แบบ ความหมายตายตัว ห้ามคิดแบบที่ห้ากลางเล่ม
%
%  colbacktitle ต้องตั้งให้เท่ากับ colback มิฉะนั้น tcolorbox จะใช้ colframe
%  (สีเข้ม) เป็นพื้นหลังของหัวกล่อง แล้วตัวอักษรหัวกล่องสีเข้มจะอ่านไม่ออก
% ---------------------------------------------------------------------
\tcbset{guidebox/.style={
  enhanced, breakable, boxrule=0pt, leftrule=2.5pt, arc=1pt,
  left=8pt, right=8pt, top=6pt, bottom=6pt,
  fonttitle=\sansall\bfseries\small, titlerule=0pt,
  toptitle=1pt, bottomtitle=1pt}}

% ฟ้า — ใจความที่ต้องจำ / ประโยคที่ร้อยหลายบทเข้าด้วยกัน
\newtcolorbox{keybox}[1][]{guidebox, colback=wash, colframe=ctr,
  colbacktitle=wash, coltitle=ctr!85!black, #1}
% ส้ม — กับดัก หรือข้อสมมติที่ถ้าลืมแล้วข้อสรุปจะผิด
\newtcolorbox{warnbox}[1][]{guidebox, colback=warnwash, colframe=beo,
  colbacktitle=warnwash, coltitle=beo!80!black, #1}
% เขียว — สิ่งที่ตรวจสอบกับของจริงในรีโปแล้ว
\newtcolorbox{goodbox}[1][]{guidebox, colback=goodwash, colframe=feg,
  colbacktitle=goodwash, coltitle=feg!65!black, #1}
% แดง — วิธีที่ผิดจริง ๆ อย่าทำ
\newtcolorbox{badbox}[1][]{guidebox, colback=badwash, colframe=bad,
  colbacktitle=badwash, coltitle=bad, #1}

% ---------------------------------------------------------------------
%  SEMANTIC MARKUP
% ---------------------------------------------------------------------
% ศัพท์อังกฤษที่แทรกในประโยคภาษาไทย ฟอนต์เนื้อความครอบคลุมละตินอยู่แล้วจึงไม่ต้อง
% สลับฟอนต์ — เก็บแมโครไว้เป็น "เครื่องหมายเชิงความหมาย" เผื่อวันหลังอยากเปลี่ยน
% การแสดงผลทั้งเล่มพร้อมกัน
\newcommand{\en}[1]{#1}
% ชื่อไฟล์ ฟังก์ชัน หรือค่าที่ปรากฏในโค้ดจริง
% ไม่บังคับขนาดตัวอักษร ปล่อยให้รับขนาดจากบริบท (สำคัญในตารางที่ใช้ \footnotesize
% มิฉะนั้นโค้ดจะกลายเป็นตัวใหญ่กว่าข้อความรอบข้างแล้วดันคอลัมน์ล้น)
% ฟอนต์ mono ถูกย่อไว้ 0.84 แล้วใน \setmonofont จึงพอดีกับเนื้อความอยู่แล้ว
\newcommand{\code}[1]{{\ttfamily #1}}

% ชื่อไฟล์ที่มีอักษรไทยอยู่ในชื่อ — ห้ามใช้ \code เพราะ DejaVu Sans Mono ไม่มี
% สัญลักษณ์ไทย ตัวอักษรจะหายไปจาก PDF โดยขึ้นแค่ warning "Missing character"
% ใช้ Garuda (sans ที่มีทั้งไทยและละติน) แทน
\newcommand{\fnth}[1]{{\sffamily\small #1}}

% ป้ายชื่อแนวคิด — สีเดียวกับที่ใช้ในไดอะแกรมและตาราง
\newcommand{\stg}[2]{\textcolor{#1}{\textbf{\en{#2}}}}

% ---------------------------------------------------------------------
%  ไวยากรณ์ของไดอะแกรม — นิยามครั้งเดียว ทุกรูปจะดูเป็นชุดเดียวกัน
%  หมายเหตุ: node ที่มีข้อความไทยต้องเผื่อ inner ysep มากกว่าไดอะแกรมภาษาละติน
%  เพราะสระบน/ล่างกินที่แนวตั้ง
% ---------------------------------------------------------------------
\tikzset{
  blk/.style={draw=rule, line width=0.6pt, rounded corners=2.5pt,
              fill=wash, minimum height=9mm, inner xsep=6pt, inner ysep=5pt,
              align=center, font=\sansall\footnotesize},
  stage/.style={blk, fill=ctr!8, draw=ctr!45},
  bestage/.style={blk, fill=beo!8, draw=beo!45},
  festage/.style={blk, fill=feg!8, draw=feg!45},
  opstage/.style={blk, fill=ops!8, draw=ops!45},
  ar/.style={-{Stealth[length=2.2mm,width=1.8mm]}, line width=0.7pt,
             draw=muted},
  lbl/.style={font=\sansall\scriptsize, text=muted, align=center},
  tag/.style={font=\sansall\scriptsize\bfseries, align=center},
}

% ---------------------------------------------------------------------
%  ตาราง
%  คอลัมน์แบบชิดซ้าย: ข้อความจัดชิดขอบสองข้าง (justified) ในคอลัมน์แคบจะเปิด
%  ช่องไฟกว้างมาก และแย่เป็นพิเศษกับภาษาไทย เพราะ ICU ตัดไทยเป็นก้อนที่ใหญ่กว่า
%  คำภาษาอังกฤษ → ใช้ L{} แทน p{} ในทุกคอลัมน์ที่มีข้อความยาว
% ---------------------------------------------------------------------
\newcolumntype{L}[1]{>{\raggedright\arraybackslash}p{#1}}
\captionsetup{skip=6pt}

% ---------------------------------------------------------------------
%  บล็อกโค้ด
%
%  ข้อจำกัดที่ต้องรู้: listings อ่านไฟล์แบบไบต์ ไม่เข้าใจ UTF-8 หลายไบต์
%  → "ห้ามใส่ภาษาไทยในบล็อกโค้ด" คอมเมนต์ในโค้ดทุกอันเป็นภาษาอังกฤษ
%    คำอธิบายภาษาไทยให้อยู่นอกบล็อกเสมอ
% ---------------------------------------------------------------------
\lstdefinelanguage{TypeScript}{
  keywords={await, async, break, case, catch, class, const, continue, default,
            delete, do, else, enum, export, extends, false, finally, for, from,
            function, if, implements, import, in, instanceof, interface, let,
            new, null, of, private, public, readonly, return, super, switch,
            this, throw, true, try, type, typeof, undefined, var, void, while,
            yield, satisfies, as},
  keywordstyle=\color{ctr}\bfseries,
  ndkeywords={string, number, boolean, Date, Promise, z, Record, Array},
  ndkeywordstyle=\color{ops},
  sensitive=true,
  comment=[l]{//},
  morecomment=[s]{/*}{*/},
  commentstyle=\color{muted}\itshape,
  stringstyle=\color{feg!70!black},
  morestring=[b]',
  morestring=[b]",
  morestring=[b]`,
}

\lstset{
  language=TypeScript,
  basicstyle=\ttfamily\scriptsize,
  backgroundcolor=\color{codebg},
  frame=leftline,
  framerule=2pt,
  rulecolor=\color{rule},
  framesep=7pt,
  xleftmargin=10pt,
  breaklines=true,
  breakatwhitespace=true,
  columns=fullflexible,
  keepspaces=true,
  showstringspaces=false,
  aboveskip=0.9em,
  belowskip=0.9em,
  literate={->}{{$\rightarrow$}}2 {=>}{{=>}}2,
}

% ชื่อ float ภาษาไทย — ถ้าไม่ตั้ง จะได้ "Figure 3" อยู่ใต้เนื้อความภาษาไทย
\renewcommand{\figurename}{รูปที่}
\renewcommand{\tablename}{ตารางที่}
\renewcommand{\refname}{เอกสารอ้างอิง}
\renewcommand{\contentsname}{สารบัญ}
% หมายเหตุ: ห้ามใส่ \color ใน \contentsname เพราะ \tableofcontents ส่งมันผ่าน
% \MakeUppercase ซึ่งจะเปลี่ยน ctr เป็น CTR แล้ว xcolor จะฟ้องว่าไม่รู้จักสีนั้น
% ถ้าจะจัดสไตล์หัวสารบัญ ให้ทำผ่าน \titleformat{\section} แทน


% หมายเหตุ: ห้ามใช้อักขระลูกศร U+2194 ที่นี่ ฟอนต์ Garuda ไม่มีสัญลักษณ์นั้น
% ตัวอักษรจะหายไปเงียบ ๆ เหลือแค่ warning "Missing character" ในล็อก
\renewcommand{\docshort}{ระยะห่างทีมกับ spike}
\hypersetup{pdftitle={ระยะห่างระหว่างโค้ดจริงกับ spike — ULMs},
            pdfauthor={ทีมพัฒนา ULMs}}

\graphicspath{{figures/}{../spike-option-b/browser-check/shots/}}

% ป้ายชื่อสี่ชั้นของระบบ — สีเดียวกันนี้ใช้ในไดอะแกรมและตารางทุกที่ในเล่ม
\newcommand{\sctr}{\stg{ctr}{สัญญา}}
\newcommand{\sbe}{\stg{beo}{backend}}
\newcommand{\sfe}{\stg{feg}{frontend}}
\newcommand{\sops}{\stg{ops}{เครือข่าย}}

% ป้ายระดับความรุนแรง — ใช้เหมือนกันทั้งเล่ม ห้ามคิดระดับที่ห้า
\newcommand{\pzero}{\textcolor{bad}{\textbf{\en{P0}}}}
\newcommand{\pone}{\textcolor{beo}{\textbf{\en{P1}}}}
\newcommand{\ptwo}{\textcolor{ctr}{\textbf{\en{P2}}}}
\newcommand{\pthree}{\textcolor{feg}{\textbf{\en{P3}}}}

\begin{document}

% =============================================================================
%  หน้าปก
% =============================================================================
\begin{titlepage}
\thispagestyle{empty}
\begin{tikzpicture}[remember picture, overlay]
  \fill[ctr!8] (current page.north west) rectangle
       ([yshift=-9.5cm]current page.north east);
  \fill[ctr] ([yshift=-9.5cm]current page.north west) rectangle
       ([yshift=-9.62cm]current page.north east);
\end{tikzpicture}

\vspace*{1.4cm}
{\sansall\color{muted}\small สำหรับทีมพัฒนา ULMs --- อ่านก่อนตัดสินใจว่าจะต่อ \en{frontend} เข้ากับ \en{backend} ด้วยวิธีไหน}

\vspace{0.9cm}
{\sansall\bfseries\fontsize{27}{33}\selectfont\color{ctr!90!black}
เราห่างจากจุดที่\\ต่อกันได้จริงแค่ไหน\par}

\vspace{0.5cm}
{\sansall\fontsize{13.5}{19}\selectfont\color{ink}
เดินตามเส้นทางข้อมูล 5 สถานี · ทุกจุดที่ \en{spike} แก้จากโค้ดจริง\\
และจุดผิดพลาดที่ต้องแก้ก่อน เรียงตามความรุนแรง\par}

\vspace{2.6cm}

\begin{minipage}{0.68\textwidth}
วันที่ 10 สิงหาคม 2569 มีการทดลองเอา \code{frontend/} ของทีมทั้งก้อน
ไปต่อกับ \en{backend} ที่เขียนขึ้นจริง แล้ว\textbf{ล็อกอินสำเร็จทั้งสามบทบาท}
ผ่านการตรวจ 13 ข้อที่ชั้นข้อมูล และ 15 ข้อในเบราว์เซอร์จริง

\vspace{0.6em}
เอกสารนี้ไม่ได้เขียนเพื่อฉลองว่าสำเร็จ แต่เขียนเพื่อตอบว่า
\textbf{ระหว่างทางต้องแก้อะไรบ้าง} และ \textbf{อะไรที่ยังพังอยู่ในโค้ดปัจจุบัน
โดยที่ยังไม่มีใครเห็น} เพราะข้อมูลจำลองปิดบังไว้

\vspace{0.6em}
จุดสำคัญที่สุดไม่ได้อยู่ที่ \en{frontend} หรือ \en{backend}
แต่อยู่ที่\textbf{ฐานข้อมูล} --- ภาค 6
\end{minipage}

\vfill
\begin{tikzpicture}
  \draw[rule, line width=0.6pt] (0,0) -- (\textwidth,0);
\end{tikzpicture}
\vspace{0.4em}

{\sansall\footnotesize\color{muted}
เอกสารฉบับที่ 3 ของชุดสัญญา \en{tRPC} · ฉบับที่ 1 \code{trpc-guide.tex} หลักการ ·
ฉบับที่ 2 \code{trpc-meeting.tex} วาระประชุม\\
ปรับปรุง 10 สิงหาคม 2569 · ข้อเท็จจริงทุกข้อสืบกลับไฟล์จริงได้ตามภาคผนวก ง}
\end{titlepage}

% =============================================================================
%  วิธีอ่านเอกสารนี้
% =============================================================================
\section*{วิธีอ่านเอกสารนี้}
\addcontentsline{toc}{section}{วิธีอ่านเอกสารนี้}

เอกสารนี้\textbf{เดินตามเส้นทางที่ข้อมูลเดินจริง} ไม่ได้เรียงตามหัวข้อทฤษฎี
เริ่มจากนิ้วที่กดปุ่มบนหน้าจอ แล้วไล่ตามไปจนถึงแถวในฐานข้อมูล
ผ่านทั้งหมด 5 สถานี

เหตุผลที่เรียงแบบนี้: ปัญหาที่เจอเวลาต่อ \en{frontend} กับ \en{backend}
มักไม่ได้อยู่ที่ปลายทางใดปลายทางหนึ่ง แต่อยู่ที่\textbf{รอยต่อ}
ซึ่งจะมองไม่เห็นเลยถ้าอ่านแยกเป็น ``บทของ \en{frontend}'' กับ ``บทของ \en{backend}''

ทุกสถานีตอบสามคำถามเดิมเสมอ

\begin{enumerate}
  \item \textbf{ทีมมีอะไรตอนนี้} --- อ้างไฟล์และบรรทัดจริงในรีโป
  \item \textbf{\en{spike} เปลี่ยนอะไร} --- โค้ดก่อน/หลังของจริง
  \item \textbf{ยังพังตรงไหน} --- สิ่งที่ \en{spike} ไม่ได้แก้ พร้อมระดับความรุนแรง
\end{enumerate}

\subsection*{ระดับความรุนแรง}

\begin{center}
\footnotesize\sansall
\begin{tabular}{@{}lL{11.6cm}@{}}
\toprule
\pzero & แก้ก่อนใครทำงานต่อได้ ถ้าปล่อยไว้ งานที่ทำเพิ่มจะต้องรื้อ หรือระบบไม่ปลอดภัย \\
\pone & แก้ก่อนเริ่มลงมือทำหน้าที่มีฟอร์ม ต้นทุนโตขึ้นทุกสัปดาห์ที่เลื่อน \\
\ptwo & แก้ตอนแตะโค้ดส่วนนั้นพอดี ไม่ต้องหยุดงานอื่นเพื่อมาแก้ \\
\pthree & งานเก็บกวาด ทำตอนไหนก็ได้ \\
\bottomrule
\end{tabular}
\end{center}

\subsection*{กล่องสี่แบบ}

\begin{center}
\footnotesize\sansall
\begin{tabular}{@{}lL{11.6cm}@{}}
\toprule
\textcolor{ctr}{\textbf{กล่องฟ้า}} & ใจความที่ต้องจำ ประโยคที่ร้อยหลายภาคเข้าด้วยกัน \\
\textcolor{beo}{\textbf{กล่องส้ม}} & กับดัก หรือข้อสมมติที่ถ้าลืมแล้วข้อสรุปจะผิด \\
\textcolor{feg}{\textbf{กล่องเขียว}} & สิ่งที่ทดสอบกับของจริงแล้ว มีผลการรันยืนยัน \\
\textcolor{bad}{\textbf{กล่องแดง}} & วิธีที่ผิดจริง ๆ อยู่ในโค้ดตอนนี้และต้องเอาออก \\
\bottomrule
\end{tabular}
\end{center}

\subsection*{สีประจำชั้นของระบบ}

สีเดียวกันนี้ใช้ทั้งในเนื้อความ ไดอะแกรม และตารางตลอดเล่ม
เห็นสีแล้วควรรู้ทันทีว่ากำลังพูดถึงชั้นไหน

\begin{center}
\footnotesize\sansall
\begin{tabular}{@{}lL{11.0cm}@{}}
\toprule
\sfe & สิ่งที่รันในเบราว์เซอร์ผู้ใช้ \\
\sctr & ข้อตกลงเรื่องรูปร่างข้อมูล สิ่งที่เดินทางระหว่างสองฝั่ง \\
\sops & เครือข่าย คุกกี้ \en{CORS} สิ่งที่อยู่ระหว่างทาง \\
\sbe & สิ่งที่รันบนเซิร์ฟเวอร์ รวมถึงฐานข้อมูล \\
\bottomrule
\end{tabular}
\end{center}

\begin{warnbox}[title={ข้อจำกัดของเอกสารนี้ที่ต้องรู้ก่อน}]
\en{spike} ที่อ้างถึงต่อ\textbf{เฉพาะระบบล็อกอิน}เท่านั้น ไม่ใช่ทั้งระบบ
เพราะเป็นเส้นเดียวที่ทั้งสองฝั่งมีของจริงพร้อมกัน ณ วันที่ทดลอง

ข้อสรุปเรื่อง ``ต่อกันได้'' จึงหมายถึง ``กลไกการต่อใช้ได้จริง''
ไม่ได้แปลว่าทั้งระบบพร้อมต่อ ส่วนที่เหลืออีก 15 \en{procedure} ยังไม่ได้เขียน
\end{warnbox}

\clearpage
\tableofcontents
\clearpage

% =============================================================================
\part{ภาพรวม --- ทีมอยู่ตรงไหน \en{spike} อยู่ตรงไหน}
% =============================================================================

\section{สามคำถามที่เอกสารนี้ตอบ}

คำถามที่ทำให้เกิดเอกสารนี้มีสามข้อ และทั้งสามข้อตอบด้วยข้อมูลจากรีโปจริง
ไม่ใช่การประเมิน

\begin{enumerate}
  \item \en{spike} แก้อะไรจากโค้ดจริงบ้าง --- ตอบครบทุกไฟล์ในภาคผนวก ก
  \item ทีมห่างจากจุดที่ต่อกันได้จริงแค่ไหน --- ตอบในหัวข้อ \ref{sec:distance}
  \item จุดผิดพลาดสำคัญที่ต้องแก้ก่อนคืออะไร --- ตอบตลอดเล่ม สรุปในภาค 7
\end{enumerate}

\begin{keybox}[title={ข้อสรุปสามบรรทัด ถ้าอ่านต่อไม่ไหว}]
\textbf{หนึ่ง} กลไกการต่อใช้ได้จริง พิสูจน์แล้วด้วยการล็อกอินสำเร็จทั้งสามบทบาท
ผ่านการตรวจ 28 ข้อ (13 ที่ชั้นข้อมูล + 15 ในเบราว์เซอร์)

\textbf{สอง} ราคาที่ต้องจ่ายก่อนเริ่มคือ\textbf{ยก \code{zod} จาก 3 เป็น 4}
ตอนนี้ยกแล้วไม่พังเลย แต่ต้นทุนจะโตตามจำนวนฟอร์มที่ยังไม่ได้ทำ

\textbf{สาม} ปัญหาที่หนักที่สุดไม่ได้อยู่ที่โค้ดฝั่งไหนเลย แต่อยู่ที่
\textbf{ฐานข้อมูล} --- ไม่มี \code{@unique} บนคอลัมน์ที่ใช้ล็อกอิน
และไม่มีตารางเก็บ \en{session} ทั้งสองข้อต้องเขียน \en{migration} ซึ่งแปลว่า
ต้องตกลงกันทั้งทีมก่อน
\end{keybox}

\section{แผนที่ 5 สถานี}

รูปที่~\ref{fig:map} คือโครงของทั้งเล่ม แต่ละกล่องคือหนึ่งภาค
ตัวเลขในวงกลมคือจำนวนจุดที่ต้องแก้ที่พบในสถานีนั้น

\begin{figure}[htbp]
\centering
% วางด้วยพิกัดตรง ๆ แทน relative positioning เพราะแบบ relative เคยทำให้
% กล่องคำอธิบายไปทับกล่องสถานีที่ 5 โดยไม่มี error แจ้งเตือน
% ระยะห่างระหว่างจุดศูนย์กลาง 3.25 ซม. กล่องกว้าง 2.85 ซม. รวม 15.85 ซม.
% ซึ่งพอดีกับ \textwidth 16.4 ซม.
\begin{tikzpicture}[
  st/.style={minimum width=2.85cm, minimum height=1.6cm, font=\sansall\scriptsize},
  cnt/.style={circle, fill=#1, text=white, font=\sansall\bfseries\scriptsize,
              inner sep=1.7pt}]

% --- ห้าสถานีเรียงเป็นเส้นเดียว ตามที่ข้อมูลเดินจริง
\node[festage, st] (s1) at (0,0)    {สถานีที่ 1\\[1pt]หน้าจอ + ฟอร์ม};
\node[stage,   st] (s2) at (3.25,0) {สถานีที่ 2\\[1pt]สัญญา + \en{type}};
\node[opstage, st] (s3) at (6.5,0)  {สถานีที่ 3\\[1pt]เครือข่าย\\[1pt]+ \en{session}};
\node[bestage, st] (s4) at (9.75,0) {สถานีที่ 4\\[1pt]\en{server}};
\node[blk, fill=bad!8, draw=bad!45, st] (s5) at (13.0,0) {สถานีที่ 5\\[1pt]ฐานข้อมูล};

\draw[ar] (s1) -- (s2);
\draw[ar] (s2) -- (s3);
\draw[ar] (s3) -- (s4);
\draw[ar] (s4) -- (s5);

% --- ป้ายภาคอยู่เหนือกล่อง
\node[lbl, above=1.6mm of s1] {ภาค 2};
\node[lbl, above=1.6mm of s2] {ภาค 3};
\node[lbl, above=1.6mm of s3] {ภาค 4};
\node[lbl, above=1.6mm of s4] {ภาค 5};
\node[lbl, above=1.6mm of s5] {ภาค 6};

% --- จำนวนจุดที่ต้องแก้ ติดมุมบนซ้ายของแต่ละกล่อง
\node[cnt=feg] at ([xshift=-1.2cm,yshift=0.6cm]s1.center) {5};
\node[cnt=ctr] at ([xshift=-1.2cm,yshift=0.6cm]s2.center) {4};
\node[cnt=ops] at ([xshift=-1.2cm,yshift=0.6cm]s3.center) {3};
\node[cnt=beo] at ([xshift=-1.2cm,yshift=0.6cm]s4.center) {1};
\node[cnt=bad] at ([xshift=-1.2cm,yshift=0.6cm]s5.center) {6};

% --- คำอธิบายอยู่ใต้ทุกกล่อง เต็มความกว้าง ไม่ทับใคร
\node[blk, fill=white, draw=rule, align=left, font=\sansall\scriptsize,
      text width=14.6cm] at (6.5,-2.3)
  {\textbf{ตัวเลขในวงกลม} = จำนวนจุดที่ต้องแก้ที่พบในสถานีนั้น รวมทั้งหมด 19 จุด\\[2pt]
   สถานีที่ 5 เป็นสีแดงเพราะเป็นสถานีเดียวที่แก้แล้วกระทบทุกคนในทีมพร้อมกัน
   --- ทุกจุดในนั้นต้องเขียน \en{migration} หรือต้องตกลงกันก่อน};

\end{tikzpicture}
\caption{เส้นทางที่ข้อมูลเดินจริงเมื่อผู้ใช้กดปุ่มเข้าสู่ระบบ และภาคที่อธิบายแต่ละสถานี}
\label{fig:map}
\end{figure}

\section{ระยะห่างระหว่างทีมกับ \en{spike}}
\label{sec:distance}

ตารางที่~\ref{tab:distance} คือคำตอบตรง ๆ ของคำถาม ``ห่างแค่ไหน''
ตัวเลขทุกช่องนับจากไฟล์จริงในรีโป

\begin{table}[htbp]
\centering
\footnotesize\sansall
\caption{สภาพของแต่ละส่วน ณ วันที่ 10 สิงหาคม 2569}
\label{tab:distance}
\begin{tabular}{@{}L{3.2cm}L{4.6cm}L{4.6cm}@{}}
\toprule
\textbf{ส่วน} & \textbf{ทีมมีอะไรตอนนี้} & \textbf{\en{spike} ต้องเพิ่มอะไร} \\
\midrule
\sfe\ โครง &
โครงเสร็จหมด \en{routing}, \en{RBAC}, \en{AppShell}, \en{ui-kit}, \en{i18n}, \en{theme} &
ไม่ต้องแตะเลย --- สวมกับสัญญาได้พอดี \\
\addlinespace[2pt]
\sfe\ หน้าจอ &
เสร็จจริง 8 จาก 26 หน้า (\en{login} + \en{admin} 6 หน้า + 1) &
ไม่เกี่ยว --- \en{spike} ต่อเฉพาะ \en{auth} \\
\addlinespace[2pt]
\sfe\ ข้อมูล &
จำลองทั้งหมด &
เปลี่ยน 3 เส้นเป็นของจริง (\en{login}/\en{me}/\en{logout}) \\
\addlinespace[2pt]
\sctr\ สัญญา &
\textbf{ยังไม่มี} &
ไฟล์ที่ \en{generate} มา 56 บรรทัด \\
\addlinespace[2pt]
\sctr\ \code{zod} &
\code{\textasciicircum 3.23.8} &
\textbf{ต้องยกเป็น 4.4.3} \\
\addlinespace[2pt]
\sops\ คุกกี้ &
\code{api-client.ts} เตรียม \code{credentials: include} ไว้แล้ว &
ต้องตั้ง \en{CORS} ให้ระบุ \en{origin} ตรง ๆ \\
\addlinespace[2pt]
\sbe\ \en{auth} &
\textbf{ไม่มีเลย} --- \code{src/} มีแต่โครง \en{Nest} เปล่ากับ \en{Prisma} &
8 ไฟล์ใหม่ \\
\addlinespace[2pt]
\sbe\ ฐานข้อมูล &
33 \en{model} 1 \en{migration} &
\textbf{ต้องมี \en{migration} เพิ่ม} ---
\en{spike} เลี่ยงไปก่อนโดยไม่แก้ \\
\bottomrule
\end{tabular}
\end{table}

อ่านตารางนี้แล้วภาพที่ควรได้คือ: \textbf{ฝั่ง frontend ใกล้มาก ฝั่ง backend ไกลมาก
และฐานข้อมูลคือที่ที่ยังไม่มีใครแตะ}

\begin{keybox}[title={ตัวเลขที่ควรจำจากตารางนี้}]
\en{frontend} ของทีมต้องแก้แค่ \textbf{14 ไฟล์ เพิ่ม 4 ลบ 1}
ซึ่งน้อยกว่าที่ทุกคนกลัวกันมาก

แต่ \en{backend} ต้องเขียนใหม่ \textbf{8 ไฟล์} จากศูนย์
และฐานข้อมูลต้องมี \en{migration} ที่ยังไม่มีใครเขียน
\end{keybox}

\section{สิ่งที่พิสูจน์แล้วจริง ๆ}

\begin{goodbox}[title={ผลการทดสอบ 28 ข้อ ผ่านหมด}]
\textbf{ชั้นข้อมูล 13 ข้อ} (\code{spike-option-b/frontend-real/e2e-check.mjs})
ใช้ \code{@trpc/client} ตัวเดียวกับที่แอปใช้ --- ยังไม่ล็อกอินแล้วถาม \code{auth.me},
ล็อกอินด้วยบัญชีภายใน, คุกกี้ติดไปกับคำขอถัดไป, รหัสผ่านผิด, อีเมลไม่ใช่ \en{KU},
ข้อความยาวเกิน \code{max(200)}, อาจารย์ล็อกอินด้วยอีเมล \en{KU},
นิสิตเครดิตต่ำได้ \en{D3}, \en{logout} แล้วคุกกี้หาย

\textbf{ชั้นหน้าจอ 15 ข้อ} (\code{spike-option-b/browser-check/ui-check.mjs})
เปิด \en{Firefox} จริง กรอกฟอร์มจริง --- รวมสามข้อที่ทดสอบด้วยวิธีอื่นไม่ได้เลย คือ
คุกกี้เป็น \code{httpOnly} จริง (\code{document.cookie} อ่านไม่เจอ),
กด \en{F5} แล้ว \en{session} ไม่หลุด, และ \en{RBAC} ยึดจริงเมื่อพิมพ์ \en{URL} ตรง ๆ
\end{goodbox}

รูปที่~\ref{fig:shell} คือภาพหน้าจอจริงหลังล็อกอินสำเร็จด้วยบัญชี \code{test\_staff}
สังเกตว่าเนื้อหน้าเขียนว่า \en{``Under development''} --- ถูกต้องแล้ว
เพราะหน้านั้นยังเป็น \en{placeholder} สิ่งที่รูปนี้พิสูจน์คือ\textbf{กรอบรอบ ๆ}
ทั้งหมดมาจากเซิร์ฟเวอร์จริง ไม่ใช่ข้อมูลจำลอง

\begin{figure}[htbp]
\centering
\includegraphics[width=0.92\textwidth]{03-staff-shell.png}
\caption{หน้าจอหลังล็อกอินด้วยบัญชีจริงจากฐานข้อมูล เมนูฝั่งซ้ายมาจากบทบาทที่เซิร์ฟเวอร์ตอบกลับ ไม่ใช่บทบาทที่เบราว์เซอร์เลือกเอง}
\label{fig:shell}
\end{figure}

\clearpage

% =============================================================================
\part{สถานีที่ 1 --- หน้าจอและฟอร์ม}
% =============================================================================

สถานีแรกคือสิ่งที่ผู้ใช้เห็นและแตะ ข่าวดีของสถานีนี้คือ\textbf{โครงทั้งหมดใช้ได้เลย}
ข่าวร้ายคือ ตรรกะการตัดสินใจว่า ``ใครเป็นใคร'' ถูกเขียนไว้ผิดที่

\section{ทีมมีอะไรตอนนี้}

โครงพื้นฐานเสร็จหมดและไม่ต้องแตะเลยตอนต่อ \en{backend} ได้แก่ \en{routing} พร้อม
\en{RBAC guard}, \en{AppShell} (แถบข้าง/แถบบน), ชุด \en{ui-kit}, \en{i18n}
ไทย--อังกฤษ, สลับธีมสว่าง/มืด และ \en{design tokens}

ส่วนเนื้อหน้าจอเสร็จไม่เท่ากันมาก ตามตารางที่~\ref{tab:pages}

\begin{table}[htbp]
\centering
\footnotesize\sansall
\caption{สถานะหน้าจอทั้ง 26 หน้า นับจากไฟล์จริงใน \code{frontend/src/features/}}
\label{tab:pages}
\begin{tabular}{@{}L{2.6cm}cL{8.4cm}@{}}
\toprule
\textbf{กลุ่ม} & \textbf{เสร็จ} & \textbf{รายละเอียด} \\
\midrule
\en{auth} & 1/1 & \code{login-page.tsx} 124 บรรทัด สองวิธีเข้าระบบ \\
\en{admin} & 6/7 & \en{dashboard}, \en{users} (511 บรรทัด), \en{status}, \en{audit},
\en{config}, \en{reports} --- ขาด \en{settings} \\
\en{borrower} & 0/7 & \en{placeholder} 6 บรรทัดทั้งหมด \\
\en{staff} & 0/7 & \en{placeholder} 6 บรรทัดทั้งหมด \\
\en{supervisor} & 0/2 & \en{placeholder} \\
\en{reports} & 0/2 & \en{placeholder} \\
\midrule
\textbf{รวม} & \textbf{8/26} & หน้าที่เสร็จทุกหน้ากินข้อมูลจาก \code{features/admin/mock-data.ts} \\
\bottomrule
\end{tabular}
\end{table}

\begin{warnbox}[title={ตัวเลข 8/26 อ่านผิดได้ง่าย}]
อย่าอ่านว่า ``เสร็จ 31\%'' เพราะงานที่เหลือไม่ได้เท่ากันทุกหน้า

หน้าที่ยังไม่ทำคือ\textbf{หน้าที่มีฟอร์ม} (ยืม คืน ตรวจสภาพ อุทธรณ์)
ซึ่งเป็นหน้าที่ผูกกับ \code{zod} มากที่สุด นี่คือเหตุผลที่หัวข้อ~\ref{sec:zod}
ยืนยันว่าต้องยก \code{zod} \textbf{ตอนนี้} ไม่ใช่ตอนหน้าฟอร์มเสร็จหมดแล้ว
\end{warnbox}

\section{จุดที่ 1 --- หน้าล็อกอินเดาบทบาทเอง}

นี่คือจุดที่ร้ายแรงที่สุดในฝั่ง \en{frontend} และเป็นจุดที่ข้อมูลจำลองปิดบังไว้ได้สนิท

\begin{badbox}[title={โค้ดที่อยู่ในรีโปตอนนี้ · \code{login-page.tsx:40--43}}]
\begin{lstlisting}
// KU email -> borrower only (students & faculty).
const handleKuLogin = () => {
  loginAs("borrower");
  navigate(HOME_ROUTE_BY_ROLE.borrower, { replace: true });
};
\end{lstlisting}
\vspace{-0.4em}
คอมเมนต์เขียนไว้เองว่า ``\en{students \& faculty}'' แต่โค้ดใต้คอมเมนต์
กำหนดบทบาทเป็น \code{borrower} ตายตัว
\end{badbox}

ปัญหาคือ\textbf{อาจารย์ทุกคนก็ใช้อีเมล \en{KU}} เมื่อทดสอบกับฐานข้อมูลจริง
บัญชี \code{orawan.p@ku.th} มี \code{RoleName} เป็น \code{Professor}
ซึ่งแปลงเป็นบทบาท \code{supervisor} ไม่ใช่ \code{borrower}

โค้ดเดิมจึงจะจับอาจารย์ทั้งคณะยัดเข้าหน้าผู้ยืม พร้อมสิทธิ์ที่ผิด
และไม่มีอะไรจับได้ เพราะข้อมูลจำลองไม่เคยมีอาจารย์ที่ล็อกอินด้วยอีเมล \en{KU}

รูปที่~\ref{fig:whodecides} อธิบายว่าความผิดพลาดนี้เกิดจากการวาง
``ผู้ตัดสินใจ'' ไว้ผิดฝั่ง

\begin{figure}[htbp]
\centering
\begin{tikzpicture}

% ===== แถวบน: วิธีเดิม
\node[lbl, anchor=west] at (0,2.75) {\textbf{วิธีเดิม} --- เบราว์เซอร์ตัดสินใจเอง};
\node[festage, minimum width=2.5cm] (a1) at (1.4,1.9) {ฟอร์ม};
\node[festage, minimum width=3.0cm, fill=bad!12, draw=bad!50] (a2) at (5.3,1.9)
  {\en{loginAs("borrower")}};
\node[festage, minimum width=2.5cm] (a3) at (9.2,1.9) {หน้าผู้ยืม};
\draw[ar] (a1) -- (a2);
\draw[ar] (a2) -- (a3);
\node[lbl, text=bad, anchor=west] at (11.0,1.9) {เซิร์ฟเวอร์\\ไม่เคยถูกถาม};

% ===== แถวล่าง: วิธีใหม่
\node[lbl, anchor=west] at (0,0.75) {\textbf{วิธีที่ \en{spike} ใช้} --- เซิร์ฟเวอร์ตัดสินใจ};
\node[festage, minimum width=2.5cm] (b1) at (1.4,-0.1) {ฟอร์ม};
\node[bestage, minimum width=3.0cm] (b2) at (5.3,-0.1) {\en{auth.login}};
\node[festage, minimum width=2.5cm] (b3) at (9.2,-0.1) {หน้าตามบทบาท};
\draw[ar] (b1) -- node[lbl, above] {\scriptsize อีเมล+รหัส} (b2);
\draw[ar] (b2) -- node[lbl, above] {\scriptsize \en{role}} (b3);
\node[lbl, text=feg, anchor=west] at (11.0,-0.1) {อ่าน \en{RoleInfo}\\จากฐานข้อมูล};

\end{tikzpicture}
\caption{ความต่างไม่ได้อยู่ที่จำนวนบรรทัด แต่อยู่ที่ว่าใครเป็นคนตอบคำถาม ``คนนี้เป็นใคร''}
\label{fig:whodecides}
\end{figure}

\begin{goodbox}[title={\en{spike} แก้เป็นอย่างนี้ · ทดสอบผ่านในเบราว์เซอร์จริงแล้ว}]
\begin{lstlisting}
const submit = async (method, identifier, password) => {
  const user = await login.mutateAsync({ method, identifier, password });
  navigate(HOME_ROUTE_BY_ROLE[user.role], { replace: true });
  return null;
};
\end{lstlisting}
\vspace{-0.4em}
ปลายทางมาจาก \code{user.role} ที่เซิร์ฟเวอร์ตอบ ไม่ใช่ค่าที่หน้าเว็บรู้ล่วงหน้า
ผลการทดสอบ: \code{orawan.p@ku.th} ไปโผล่ที่ \code{/supervisor/approvals} ถูกต้อง
\end{goodbox}

\section{จุดที่ 2 --- บทบาทเก็บอยู่ใน \en{localStorage}}

จุดนี้รุนแรงกว่าจุดแรก เพราะไม่ใช่แค่ ``เดาผิด'' แต่เป็น ``ให้ผู้ใช้เลือกเองได้''

\begin{badbox}[title={\pzero\ ช่องโหว่ · \code{auth.store.ts:36, 46}}]
\begin{lstlisting}
// login
window.localStorage.setItem(MOCK_USER_STORAGE_KEY, role);
// on app start
role = window.localStorage.getItem(MOCK_USER_STORAGE_KEY);
set({ user: isRole(role) ? MOCK_USERS[role] : null, isLoading: false });
\end{lstlisting}
\vspace{-0.4em}
\en{localStorage} แก้ได้จาก \en{DevTools} ของเบราว์เซอร์ในสามวินาที
ใครก็ตามที่เปิดหน้าเว็บนี้ พิมพ์คำเดียวก็เป็น \code{admin} ได้
\end{badbox}

ตอนนี้ยังไม่อันตรายจริง เพราะหลังบ้านยังไม่มีข้อมูลจริงให้เข้าถึง
แต่\textbf{วันที่ backend ต่อเสร็จ ช่องนี้จะกลายเป็นประตูหลังทันที}
ถ้าไม่มีใครจำได้ว่าต้องเอาออก

\en{spike} แก้โดยลบทั้งสองฟังก์ชันทิ้ง แล้วให้เซิร์ฟเวอร์เป็นผู้ตอบแทน
ผ่านคุกกี้ \code{httpOnly} ซึ่ง \en{JavaScript} อ่านไม่ได้เลย ---
รายละเอียดกลไกอยู่ในภาค 4

\begin{keybox}[title={ประโยคที่ร้อยจุดที่ 1 กับจุดที่ 2 เข้าด้วยกัน}]
ทั้งสองจุดคือความผิดพลาดชนิดเดียวกัน: \textbf{ให้เบราว์เซอร์เป็นผู้ตัดสินใจ
เรื่องที่เบราว์เซอร์ไม่มีสิทธิ์ตัดสิน}

หลังแก้แล้ว ไม่มีอะไรในเบราว์เซอร์แจกสิทธิ์ให้ตัวเองได้อีก
นี่คือประโยชน์ที่จับต้องได้ที่สุดของการต่อ \en{backend} จริง
มากกว่าเรื่อง ``ได้ข้อมูลจริงมาแสดง''
\end{keybox}

\section{จุดที่ 3 --- รหัสผ่านทดสอบอยู่ในรีโป}

\code{frontend/src/features/auth/mock-auth.ts:62--70} เก็บตารางบัญชีและรหัสผ่าน
เป็นข้อความธรรมดา และ \code{frontend/README.md:14--18} พิมพ์ซ้ำอีกรอบเป็นตาราง

ทั้งสองที่มีคำเตือนเขียนกำกับไว้แล้วว่าให้ลบเมื่อ \en{backend} มา
ซึ่งดีมาก แต่\textbf{คำเตือนไม่ใช่กลไก} --- ไม่มีอะไรบังคับให้ลบจริง

\en{spike} ลบ \code{mock-auth.ts} ทิ้งทั้งไฟล์ (70 บรรทัด) บัญชีทดสอบย้ายไปอยู่ใน
\code{seed.ts} ของฐานข้อมูล ซึ่งเป็นที่ที่ถูกต้อง เพราะรหัสผ่านถูก \en{hash}
ด้วย \en{scrypt} ก่อนเก็บ ไม่ได้อยู่เป็นข้อความธรรมดาอีกต่อไป

\section{จุดที่ 4 --- ฟอร์มต้องกลายเป็น \en{async}}

จุดนี้ไม่ใช่ความผิดพลาด แต่เป็น\textbf{ผลกระทบลูกโซ่}ที่ทีมควรรู้ล่วงหน้า
เพราะจะเกิดซ้ำกับทุกฟอร์มที่เหลือ

ตอนที่การล็อกอินเป็นของจำลอง มันตอบทันที ฟังก์ชันจึงเขียนแบบ \en{synchronous} ได้

\begin{lstlisting}
// before: the mock answers instantly
onSubmit: (values: LocalLoginValues) => string | null | void;
// after: a real round trip
onSubmit: (values: LocalLoginValues) => Promise<string | null | void>;
\end{lstlisting}

การเปลี่ยนคำเดียวนี้ลากของตามมาอีกสามอย่างในทุกฟอร์ม

\begin{enumerate}
  \item ตัวเรียกต้อง \code{await} และตัวจัดการต้องเป็น \code{async}
  \item ปุ่มต้องมีสถานะ ``กำลังส่ง'' เพราะคำขอจริงช้าพอที่ผู้ใช้จะกดซ้ำ
  \item ต้องมีที่แสดงข้อความผิดพลาดจากเซิร์ฟเวอร์ ซึ่งฟอร์มเดิมไม่มี
\end{enumerate}

ข้อ 3 เห็นชัดที่ช่องอีเมล \en{KU}: ตอนเป็นของจำลอง เส้นทางนั้น\textbf{ล้มเหลวไม่ได้เลย}
จึงไม่เคยมีที่ให้แสดง ``รหัสผ่านไม่ถูกต้อง'' พอต่อของจริงกลายเป็นกรณีที่พบบ่อยที่สุด

\begin{warnbox}[title={คูณด้วยจำนวนฟอร์มที่ยังไม่ได้ทำ}]
งานสามข้อข้างบนใช้เวลาไม่กี่นาทีต่อฟอร์ม แต่ต้องทำ\textbf{ทุกฟอร์ม}
และหน้าที่ยังเป็น \en{placeholder} อีก 19 หน้าส่วนใหญ่มีฟอร์ม

ถ้าทีมเขียนฟอร์มที่เหลือแบบ \en{synchronous} ต่อไปโดยยังไม่ต่อ \en{backend}
งานรื้อรอบเดียวจะใหญ่กว่านี้หลายเท่า
\end{warnbox}

\section{จุดที่ 5 --- แถบข้างแสดงชื่อที่ตายตัว}

\code{sidebar.tsx} แสดง \code{persona.name} ซึ่งมาจาก \code{constants/navigation.ts}
--- เป็นค่าคงที่หนึ่งชุดต่อหนึ่งบทบาท ไม่ใช่ชื่อของคนที่ล็อกอินจริง

ตอนนี้ยังไม่มีใครเห็นความต่าง เพราะ \code{seed.ts} ตั้งชื่อให้ตรงกับข้อมูลจำลองพอดี
แต่พอดูให้ละเอียดจะพบว่าไม่ตรง: ฐานข้อมูลเก็บว่า ``อรวรรณ ภักดี''
ส่วนแถบข้างเขียนว่า ``ผศ.ดร. อรวรรณ ภักดี''

\en{spike} \textbf{ไม่ได้แก้จุดนี้} เพราะเป็นงานของเจ้าของ \en{AppShell}
ไม่ใช่ของระบบล็อกอิน จัดเป็น \ptwo\ --- แก้ตอนแตะไฟล์นั้นพอดี

\section{สรุปสถานีที่ 1}

\begin{table}[htbp]
\centering
\footnotesize\sansall
\caption{ห้าจุดของสถานีที่ 1 และสถานะ}
\label{tab:station1}
\begin{tabular}{@{}cL{6.0cm}cL{4.2cm}@{}}
\toprule
\textbf{ระดับ} & \textbf{จุด} & \textbf{\en{spike} แก้แล้ว} & \textbf{ไฟล์} \\
\midrule
\pzero & บทบาทเก็บใน \en{localStorage} & ใช่ & \code{auth.store.ts} \\
\pzero & รหัสผ่านทดสอบอยู่ในรีโป & ใช่ (ลบทั้งไฟล์) & \code{mock-auth.ts} \\
\pone & หน้าล็อกอินเดาบทบาทเอง & ใช่ & \code{login-page.tsx} \\
\ptwo & ฟอร์มต้องเป็น \en{async} & ใช่ (2 ไฟล์) & \code{login-method-*.tsx} \\
\ptwo & แถบข้างแสดงชื่อตายตัว & \textbf{ยังไม่แก้} & \code{sidebar.tsx} \\
\bottomrule
\end{tabular}
\end{table}

\clearpage

% =============================================================================
\part{สถานีที่ 2 --- สัญญาและ \en{type}}
% =============================================================================

สถานีนี้คือรอยต่อที่แท้จริง และเป็นที่ที่ค่าผ่านทางทั้งหมดของทางเลือก ข. อยู่

\section{สัญญาเดินทางมาแบบไหน}

ในทางเลือก ข. ``สัญญา'' คือ\textbf{ไฟล์หนึ่งไฟล์}ที่เครื่องมือสร้างจาก \en{router}
ของ \en{backend} แล้วคัดลอกมาวางไว้ใน \code{frontend/src/generated/server.ts}

ไฟล์นั้น \en{import} แค่สองอย่าง

\begin{lstlisting}
import { initTRPC } from "@trpc/server";
import { z } from "zod";
\end{lstlisting}

บรรทัดที่สองคือต้นเหตุของทุกอย่างในภาคนี้

\section{จุดที่ 6 --- \code{zod} คนละรุ่น}
\label{sec:zod}

\begin{badbox}[title={\pone\ ตัวบล็อกที่เจอตั้งแต่ยังไม่เขียนโค้ดบรรทัดแรก}]
\code{frontend/package.json:40} ประกาศ \code{"zod": "\textasciicircum 3.23.8"}

\code{spike-option-b/backend/package.json} ประกาศ \code{"zod": "4.4.3"}

ไฟล์สัญญาที่ \en{generate} ออกมาผูกกับรุ่นของฝั่งที่สร้างมัน คือ 4.4.3
\end{badbox}

ทำไมรุ่นต่างกันแล้วพัง --- และทำไมมันพัง\textbf{เงียบ}

\en{TypeScript} ตัดสินว่า \en{type} สองตัวเหมือนกันหรือไม่จากโครงสร้าง
แต่ \code{ZodString} ของ \code{zod@3} กับของ \code{zod@4} มีโครงสร้างภายในต่างกัน
ทั้งคู่ชื่อ \code{ZodString} เหมือนกันแต่ไม่ใช่ตัวเดียวกัน ผลคือการอ่าน \en{type}
ออกจากสัญญาจะล้มเหลว แล้วทุก \en{procedure} ตกเป็น \code{any}

\code{any} ไม่ทำให้ \en{build} พัง มันแค่แปลว่า ``ไม่ต้องตรวจอะไรเลย''
ซึ่งคือ\textbf{การสูญเสียเหตุผลทั้งหมดที่เลือกใช้ tRPC ตั้งแต่แรก}
โดยที่ไม่มีข้อความเตือนสักบรรทัด

\begin{figure}[htbp]
\centering
\begin{tikzpicture}

\node[bestage, minimum width=3.1cm] (be) at (0,0) {\en{backend}\\\scriptsize \en{zod 4.4.3}};
\node[stage, minimum width=3.1cm] (ct) at (4.6,0) {ไฟล์สัญญา\\\scriptsize อ้าง \en{ZodString} ของ 4.4.3};
\node[festage, minimum width=3.1cm] (fe) at (9.2,0) {\en{frontend}\\\scriptsize \en{zod 3.23.8}};

\draw[ar] (be) -- node[lbl, above] {\scriptsize \en{generate}} (ct);
\draw[ar] (ct) -- node[lbl, above] {\scriptsize \en{import type}} (fe);

\node[blk, fill=bad!10, draw=bad!50, text width=6.6cm, align=left,
      font=\sansall\scriptsize] (out) at (4.6,-2.3)
  {\textbf{ผลลัพธ์} \en{TypeScript} จับคู่ \en{type} ไม่ได้\\
   $\rightarrow$ ทุก \en{procedure} กลายเป็น \en{any}\\
   $\rightarrow$ \textbf{\en{build} ยังผ่าน} ไม่มีคำเตือนใด ๆ\\
   $\rightarrow$ พิมพ์ชื่อฟิลด์ผิดก็ไม่มีใครทัก};
\draw[ar, draw=bad!60] (ct) -- (out);

\end{tikzpicture}
\caption{เหตุผลที่ \code{zod} คนละรุ่นเป็นความล้มเหลวแบบเงียบ ไม่ใช่ข้อผิดพลาดที่มองเห็น}
\label{fig:zoddrift}
\end{figure}

\subsection{ข่าวดี: ยกตอนนี้ไม่พังเลย}

\begin{goodbox}[title={ทดลองแล้ว ยก \code{zod} เป็น 4.4.3 โดยยังไม่แตะโค้ดสักบรรทัด}]
ผลการรัน \code{tsc --noEmit} คือ \textbf{ผ่าน 0 ข้อผิดพลาด}

เหตุผล: ทั้งโปรเจกต์ใช้ \code{zod} แค่ \textbf{2 ไฟล์} คือ
\code{features/auth/login.schema.ts} กับ \code{schemas/loan-request.schema.ts}
และทั้งสองใช้แต่ \en{API} ที่ \code{zod@4} ยังรองรับ
(\code{.email()} และ \code{.datetime()} ถูกประกาศเลิกใช้แล้วแต่ยังทำงานได้)
\end{goodbox}

การยก \code{zod} ลากของตามมาหนึ่งอย่าง คือ \code{@hookform/resolvers}
ต้องขึ้นจาก \code{\textasciicircum 3.10.0} เป็น \code{5.2.2}
เพราะรุ่น 3 รองรับเฉพาะ \code{zod@3}

\begin{warnbox}[title={ข่าวดีนี้มีวันหมดอายุ}]
ที่ยกได้ฟรีวันนี้เพราะมี \en{schema} แค่ 2 ไฟล์

หน้าที่ยังไม่ทำอีก 19 หน้าคือหน้าที่มีฟอร์ม แต่ละฟอร์มจะเพิ่ม \en{schema}
ต้นทุนการยกจึงโตทุกสัปดาห์ที่เลื่อนออกไป

\textbf{ถ้าทีมจะเลือกทางนี้ ต้องยก \code{zod} ตอนนี้ ไม่ใช่ตอนหน้าฟอร์มเสร็จหมดแล้ว}
\end{warnbox}

\section{จุดที่ 7 --- \en{type} ของผู้ใช้ที่เขียนด้วยมือ}

\code{frontend/src/types/domain.ts:9--18} เขียน \code{interface User} ขึ้นเอง
โดยอ้างอิงความเข้าใจ ณ ตอนนั้น พอเทียบกับสิ่งที่เซิร์ฟเวอร์ส่งจริง พบว่าผิด 3 ฟิลด์

\begin{table}[htbp]
\centering
\footnotesize\sansall
\caption{สามฟิลด์ที่เดาผิด และเหตุผลจากฐานข้อมูลจริง}
\label{tab:userdrift}
\begin{tabular}{@{}L{2.7cm}L{2.4cm}L{2.4cm}L{4.9cm}@{}}
\toprule
\textbf{ฟิลด์} & \textbf{ทีมเขียนไว้} & \textbf{ของจริง} & \textbf{เพราะอะไร} \\
\midrule
\code{id} & \code{string} & \code{number} &
\code{AccountInfo.AccountKey} เป็น \code{Int @id} \\
\addlinespace[2pt]
\code{departmentId} & \code{string} & \code{string | null} &
ผู้ใช้ที่ยังไม่ถูกผูกกับภาควิชาใดจะไม่มีค่า \\
\addlinespace[2pt]
\code{maxBorrowDays} & \textbf{ไม่มีเลย} & \code{number} &
เซิร์ฟเวอร์คำนวณจาก \en{CreditTier} ส่งมาให้ \\
\bottomrule
\end{tabular}
\end{table}

ข้อ \code{departmentId} อันตรายแบบเงียบที่สุด: โค้ดที่เอาไปต่อสตริงจะได้คำว่า
\code{"null"} โผล่บนหน้าจอ ไม่ใช่ค่าว่าง

\subsection{วิธีที่ \en{spike} ใช้กันไม่ให้เกิดอีก}

เลิกเขียน \code{User} ด้วยมือ เปลี่ยนเป็นดึงออกมาจากสัญญาโดยตรง

\begin{lstlisting}
// src/types/domain.ts
export type User = ContractUser;               // from inferRouterOutputs
export type Role = ContractUser["role"];
export type CreditBand = ContractUser["creditBand"];
\end{lstlisting}

ผลทันทีเมื่อทำแบบนี้: \code{tsc} ชี้จุดที่ขัดกับเซิร์ฟเวอร์ให้เอง และมันชี้จริง ---
\code{mock-auth.ts} พัง 4 บรรทัด ซึ่งเป็นไฟล์ที่ต้องลบอยู่แล้ว

\begin{goodbox}[title={ผลที่สำคัญกว่านั้น}]
\textbf{นอกจากไฟล์ที่ต้องลบแล้ว ไม่มีอะไรพังเลยทั้งโปรเจกต์}

แปลว่ารูปร่างข้อมูลที่เซิร์ฟเวอร์ส่งมาสวมกับโค้ดที่ทีมเขียนไว้ได้พอดี
งานต่อ \en{backend} จึงไม่ใช่การรื้อ แต่เป็นการเสียบ
\end{goodbox}

\section{จุดที่ 8 --- รหัสข้อผิดพลาดคนละชุด}

\code{frontend/src/lib/error-messages.ts:11} เตรียมข้อความไว้สำหรับรหัส
\code{INVALID\_DOMAIN}

\code{backend/src/contract.types.ts} ส่งรหัส \code{NOT\_KU\_EMAIL}

\begin{badbox}[title={คนละรหัส ทั้งสองฝั่ง \en{build} ผ่าน ไม่มีใครเตือน}]
ผู้ใช้ที่กรอกอีเมล \en{gmail} จะเห็นข้อความ ``เกิดข้อผิดพลาดที่ไม่ทราบสาเหตุ''
แทนข้อความที่เขียนเตรียมไว้อย่างดี

สาเหตุเชิงโครงสร้าง: ตารางแปลรหัสประกาศเป็น \code{Record<string, string>}
ซึ่งรับ \en{key} อะไรก็ได้ การค้นหาที่ไม่เจอไม่ใช่ข้อผิดพลาด มันแค่คืนค่าว่าง
\end{badbox}

นี่คือความต่างที่ชัดที่สุดระหว่างสองทางเลือก: ในทางเลือก ก. ฝั่ง \en{frontend}
\code{import \{ BUSINESS\_ERROR \}} จากแพ็กเกจกลาง การเพิ่มรหัสใหม่แล้วไม่แปล
จะกลายเป็น \en{compile error} ทันที ในทางเลือก ข. ไม่มีอะไรเชื่อมสองฝั่งเลย

\section{จุดที่ 9 --- สัญญาเก่าค้างได้ และ \en{CI} มองไม่เห็น}

ทางเลือก ข. มีความเสี่ยงเฉพาะตัวคือ ไฟล์สัญญาเป็นสำเนา ถ้าแก้ \en{router}
แล้วลืมสั่ง \en{generate} ฝั่ง \en{frontend} จะยัง \en{build} ผ่าน
เพราะ \en{type} เก่าก็ยังเป็น \en{type} ที่ถูกต้อง

\en{spike} จึงเขียนด่าน \en{CI} ไว้สองตัวเพื่อจับเรื่องนี้ และทดสอบแล้วว่าจับได้จริง
\textbf{แต่ตอนเอา \en{frontend} จริงมาต่อ พบว่าด่านทั้งสองมองไม่เห็นมัน}

\begin{badbox}[title={\ptwo\ ทดลองทำให้พังแล้ว ด่านยังเขียวทั้งคู่}]
ตั้ง \code{frontend-real} ให้ประกาศ \code{zod} เป็น \code{\textasciicircum 4.7.0}
(ผิดทั้งรุ่นและไม่ \en{pin}) และแก้ไฟล์สัญญาให้เพี้ยนจากต้นฉบับ แล้วรันด่านทั้งสอง

\begin{lstlisting}
check-versions:       PASS
check-contract-fresh: PASS
\end{lstlisting}
\vspace{-0.4em}
สาเหตุ: สคริปต์ทั้งคู่ระบุ \en{path} ไว้ตรง ๆ ว่า \code{frontend/}
ผู้บริโภคสัญญารายที่สองจึงอยู่นอกสายตา
\end{badbox}

\begin{keybox}[title={นี่ไม่ใช่บั๊กของสคริปต์ แต่เป็นรูปร่างของปัญหา}]
ทางเลือก ข. ต้องมี\textbf{ด่านหนึ่งด่านต่อผู้บริโภคหนึ่งราย}
และต้องมีคนจำได้ว่าต้องเพิ่มด่านเมื่อมีผู้บริโภครายใหม่
(แอปมือถือ, \en{admin portal}, ชุดทดสอบ \en{e2e})

ทางเลือก ก. ไม่มีคำถามนี้ เพราะ \en{npm workspaces} ยุบ \code{zod}
ให้เหลือก้อนเดียวโดยโครงสร้าง ไม่ใช่โดยวินัย

ถ้าทีมเลือกทางนี้ ต้องแก้สคริปต์ให้\textbf{วนทุกโฟลเดอร์ที่กินสัญญา}
ไม่ใช่เขียนชื่อเดียวตายตัว
\end{keybox}

\section{สรุปสถานีที่ 2}

\begin{table}[htbp]
\centering
\footnotesize\sansall
\caption{ห้าจุดของสถานีที่ 2 และสถานะ}
\label{tab:station2}
\begin{tabular}{@{}cL{6.4cm}cL{3.8cm}@{}}
\toprule
\textbf{ระดับ} & \textbf{จุด} & \textbf{\en{spike} แก้แล้ว} & \textbf{ไฟล์} \\
\midrule
\pone & \code{zod} 3 ต้องยกเป็น 4 (ลาก \code{@hookform/resolvers} 3$\rightarrow$5 ตามมาด้วย) & ใช่ & \code{package.json} \\
\pone & \code{User} เขียนมือ ผิด 3 ฟิลด์ & ใช่ & \code{types/domain.ts} \\
\ptwo & รหัสข้อผิดพลาดคนละชุด & ใช่ & \code{error-messages.ts} \\
\ptwo & \en{CI} มองไม่เห็นผู้บริโภครายที่สอง & \textbf{ยังไม่แก้} & \code{scripts/*.mjs} \\
\bottomrule
\end{tabular}
\end{table}

\clearpage

% =============================================================================
\part{สถานีที่ 3 --- เครือข่ายและ \en{session}}
% =============================================================================

สถานีนี้คือระยะทางระหว่างเบราว์เซอร์กับเซิร์ฟเวอร์ สิ่งที่เดินทางบนเส้นนี้
ไม่ได้มีแค่ข้อมูล แต่มี\textbf{ตัวตนของผู้ใช้}ด้วย

\section{ทีมเตรียมอะไรไว้แล้ว}

ข่าวดี: \code{frontend/src/lib/api-client.ts:42} ตั้ง \code{credentials: "include"}
ไว้แล้วพร้อมคอมเมนต์ว่า ``ส่ง \en{cookie} ทุกครั้ง'' แปลว่าทีมตัดสินใจเรื่อง
รูปแบบการยืนยันตัวตนถูกต้องมาตั้งแต่ต้น --- ใช้คุกกี้ ไม่ใช่ \en{token} ใน \en{localStorage}

\code{frontend/README.md} ก็เขียนไว้ชัดว่า \en{``Auth: httpOnly cookie (session-based)''}

\begin{keybox}[title={การตัดสินใจที่ถูกแล้ว ไม่ต้องเปลี่ยน}]
เลือกคุกกี้ \code{httpOnly} แทนการเก็บ \en{token} ใน \en{localStorage}
เป็นการตัดสินใจที่ถูกและสำคัญ เพราะ \en{JavaScript} อ่านคุกกี้แบบนี้ไม่ได้เลย
สคริปต์ที่ถูกฝังในหน้าเว็บจึงขโมยไปไม่ได้

\en{spike} ทดสอบยืนยันในเบราว์เซอร์จริงแล้วว่า \code{document.cookie}
มองไม่เห็น \code{ulms\_session} จริง ๆ
\end{keybox}

\section{กลไกที่เกิดขึ้นจริงเมื่อล็อกอิน}

รูปที่~\ref{fig:seq} คือลำดับที่เกิดขึ้นจริง วัดจากการทดสอบ ไม่ใช่ภาพอุดมคติ

\begin{figure}[htbp]
\centering
\begin{tikzpicture}[font=\sansall\scriptsize]

% เส้นชีวิตสามเส้น
\node[festage, minimum width=2.6cm] (b) at (0,0) {เบราว์เซอร์};
\node[bestage, minimum width=2.6cm] (s) at (6.0,0) {เซิร์ฟเวอร์ \en{:3002}};
\node[blk, fill=bad!8, draw=bad!45, minimum width=2.6cm] (d) at (11.4,0) {ฐานข้อมูล};

\draw[rule, dashed] (b.south) -- ++(0,-6.7);
\draw[rule, dashed] (s.south) -- ++(0,-6.7);
\draw[rule, dashed] (d.south) -- ++(0,-6.7);

% ขั้นตอน
\draw[ar] (0,-1.0) -- node[lbl, above]{\en{auth.login} + \en{Origin}} (6.0,-1.0);
\draw[ar] (6.0,-1.7) -- node[lbl, above]{\en{findFirst} ตามอีเมล} (11.4,-1.7);
\draw[ar] (11.4,-2.4) -- node[lbl, above]{แถวบัญชี + \en{RoleName}} (6.0,-2.4);
\draw[ar, draw=ops] (6.0,-3.1) -- node[lbl, above, text=ops]{\textbf{\en{Set-Cookie}} \en{httpOnly}} (0,-3.1);

% เส้นคั่นบอกว่าเวลาเดินไปแล้ว ผู้ใช้กลับมาใหม่ในภายหลัง
\draw[rule, dotted] (-0.6,-3.7) -- (12.0,-3.7);
\node[lbl, anchor=west, text=muted, fill=white, inner sep=2pt] at (0.2,-3.7)
  {ผู้ใช้กด \en{F5} --- หน้าเว็บลืมทุกอย่างที่อยู่ในหน่วยความจำ};

\draw[ar] (0,-4.6) -- node[lbl, above]{\en{auth.me} + คุกกี้ที่เบราว์เซอร์แนบให้เอง} (6.0,-4.6);
\draw[ar] (6.0,-5.4) -- node[lbl, above]{อ่านบทบาท\textbf{ใหม่}ทุกครั้ง} (11.4,-5.4);
\draw[ar] (6.0,-6.2) -- node[lbl, above]{ผู้ใช้คนเดิม ยังอยู่ในระบบ} (0,-6.2);

\end{tikzpicture}
\caption{ลำดับจริงของการล็อกอินและการกู้ \en{session} คืนหลังรีเฟรช --- ทั้งสองเส้นทางถูกทดสอบในเบราว์เซอร์จริงแล้ว}
\label{fig:seq}
\end{figure}

จุดที่ควรสังเกตในรูปคือลูกศรเส้นที่ห้า: หลังรีเฟรช เบราว์เซอร์แนบคุกกี้ให้เอง
โดยที่โค้ดไม่ต้องทำอะไรเลย และเซิร์ฟเวอร์\textbf{อ่านบทบาทใหม่จากฐานข้อมูลทุกครั้ง}
ไม่ได้ฝังบทบาทไว้ในคุกกี้

การอ่านใหม่ทุกครั้งช้ากว่า แต่แลกมาด้วยความถูกต้อง: วันที่ผู้ดูแลถอนสิทธิ์ใคร
คนนั้นเสียสิทธิ์ทันที ไม่ใช่ตอนคุกกี้หมดอายุ

\section{จุดที่ 10 --- \en{CORS} ที่ยังไม่มี}

\code{backend/src/main.ts} ของทีมตอนนี้มี 7 บรรทัด และ\textbf{ไม่มี \en{CORS}
ไม่มี \en{cookie-parser}}

\begin{badbox}[title={\pone\ ถ้าไม่ตั้งสองอย่างนี้ ทุกอย่างที่เหลือใช้ไม่ได้}]
\begin{lstlisting}
// backend/src/main.ts as it stands today
const app = await NestFactory.create(AppModule);
await app.listen(process.env.PORT ?? 3000);
\end{lstlisting}
\vspace{-0.4em}
เบราว์เซอร์จะทิ้งคำตอบทุกอันก่อนที่โค้ดฝั่งหน้าเว็บจะได้เห็น
และเซิร์ฟเวอร์จะอ่านคุกกี้ที่ส่งมาไม่ได้
\end{badbox}

สิ่งที่ต้องเพิ่ม และเหตุผลของแต่ละบรรทัด

\begin{lstlisting}
app.use(cookieParser());          // without this, req.cookies is undefined
app.enableCors({
  origin: allowedOrigins,         // must be an explicit list, NOT "*"
  credentials: true,              // without this the browser drops the cookie
  methods: ['GET', 'POST', 'OPTIONS'],
});
\end{lstlisting}

\begin{warnbox}[title={กับดักที่คนเจอบ่อยที่สุดเรื่อง \en{CORS} กับคุกกี้}]
เมื่อ \code{credentials: true} มาตรฐานของเบราว์เซอร์\textbf{ห้าม}ใช้
\code{origin: "*"} ต้องระบุที่อยู่ให้ครบทุกอันแบบเจาะจง

แปลว่าทุกครั้งที่มีคนรัน \en{frontend} บนพอร์ตใหม่ ต้องไปเพิ่มชื่อในฝั่งเซิร์ฟเวอร์ด้วย
--- ในการทดลองครั้งนี้ต้องเพิ่ม \code{http://localhost:5275} เข้าไปจริง ๆ
ก่อนถึงจะล็อกอินได้
\end{warnbox}

\section{จุดที่ 11 --- \en{proxy} ที่ชี้ผิดที่}

\code{frontend/vite.config.ts} ตั้ง \en{proxy} ให้ \code{/api} วิ่งไปที่
\code{localhost:3000} ซึ่งเป็นการตั้งค่าสำหรับ \en{REST} ไม่ใช่ \en{tRPC}

ถ้าใช้ \en{tRPC} แบบที่ \en{spike} ทำ คือคุยตรงกับ \code{:3002} พร้อม \en{CORS}
\en{proxy} ตัวนี้จะไม่ถูกใช้เลย จัดเป็น \ptwo\ --- ไม่ได้พัง แต่ทำให้สับสน
และควรตัดสินใจให้ตรงกันว่าจะใช้แบบไหน

\section{จุดที่ 12 --- แยก ``ต่อไม่ติด'' ออกจาก ``รหัสผิด'' ไม่ได้}

\code{error-messages.ts} เดิมไม่มีทางแยกสองกรณีนี้ ผลคือ
\textbf{ถ้าเซิร์ฟเวอร์ดับ ผู้ใช้จะเห็นข้อความว่ารหัสผ่านไม่ถูกต้อง} ซึ่งเป็นการโกหก
และทำให้ผู้ใช้พิมพ์รหัสซ้ำอีกสิบรอบโดยไม่มีทางสำเร็จ

\en{spike} แก้โดยดูว่าคำตอบมีข้อมูลจากเซิร์ฟเวอร์แนบมาหรือไม่

\begin{lstlisting}
export function isConnectionError(error: unknown): boolean {
  // tRPC still throws TRPCClientError on transport failure, but with no
  // server payload attached. That absence is the signal.
  return error instanceof TRPCClientError && !error.data;
}
\end{lstlisting}

\begin{goodbox}[title={ทดลองได้เองใน 30 วินาที}]
ปิด \en{backend} แล้วลองล็อกอิน ต้องขึ้นข้อความว่า
``เชื่อมต่อเซิร์ฟเวอร์ไม่ได้ กรุณาตรวจสอบว่า \en{backend} ทำงานอยู่''
ไม่ใช่ ``ชื่อผู้ใช้หรือรหัสผ่านไม่ถูกต้อง''
\end{goodbox}

\section{สรุปสถานีที่ 3}

\begin{table}[htbp]
\centering
\footnotesize\sansall
\caption{สามจุดของสถานีที่ 3 และสถานะ}
\label{tab:station3}
\begin{tabular}{@{}cL{6.4cm}cL{3.8cm}@{}}
\toprule
\textbf{ระดับ} & \textbf{จุด} & \textbf{\en{spike} แก้แล้ว} & \textbf{ไฟล์} \\
\midrule
\pone & \en{main.ts} ไม่มี \en{CORS} / \en{cookie-parser} & ใช่ (ในฝั่ง \en{spike}) & \code{backend/src/main.ts} \\
\ptwo & แยกกรณีต่อไม่ติดไม่ได้ & ใช่ & \code{lib/trpc.ts} \\
\ptwo & \en{proxy} \code{/api} ชี้ \code{:3000} & \textbf{ยังไม่ตัดสินใจ} & \code{vite.config.ts} \\
\bottomrule
\end{tabular}
\end{table}

\clearpage

% =============================================================================
\part{สถานีที่ 4 --- \en{server}}
% =============================================================================

\section{ทีมมีอะไรตอนนี้}

คำตอบสั้น ๆ คือ \textbf{ยังไม่มีระบบล็อกอินเลย} ไม่ใช่ ``มีแต่ยังไม่เสร็จ''
แต่คือยังไม่ได้เริ่ม

\code{backend/src/} ที่ไม่ใช่ไฟล์ที่เครื่องมือสร้างให้ มีอยู่ 9 ไฟล์
และเป็นโครงเปล่าของ \en{NestJS} กับการต่อ \en{Prisma} ทั้งหมด

\begin{table}[htbp]
\centering
\footnotesize\sansall
\caption{ไฟล์ทั้งหมดใน \code{backend/src/} ที่คนเขียนเอง}
\label{tab:beteam}
\begin{tabular}{@{}L{4.4cm}L{8.6cm}@{}}
\toprule
\textbf{ไฟล์} & \textbf{คืออะไร} \\
\midrule
\code{main.ts} & 7 บรรทัด ไม่มี \en{CORS} ไม่มี \en{cookie-parser} \\
\code{app.module.ts} & โครงเปล่า ยังไม่ได้ตั้ง \code{TRPCModule} \\
\code{app.controller.ts} & \en{controller} ตัวอย่างที่ \en{Nest} สร้างให้ตอน \code{new} \\
\code{app.service.ts} & \en{service} ตัวอย่าง \\
\code{prisma.service.ts} & ต่อฐานข้อมูล --- ส่วนนี้ใช้ได้จริง \\
\code{prisma.module.ts} & \en{module} ของข้างบน \\
\code{querytest.ts} & สคริปต์ลองคิวรี \\
\code{*.spec.ts} & ไฟล์ทดสอบตัวอย่าง 2 ไฟล์ \\
\bottomrule
\end{tabular}
\end{table}

\section{จุดที่ 13 --- แปดไฟล์ที่ต้องเขียนเพิ่ม}

\en{spike} เขียนไฟล์เหล่านี้ขึ้นมาแล้วและใช้งานได้จริง ยกมาได้เลย
รายการนี้คือ ``ระยะห่าง'' ฝั่ง \en{backend} ที่เป็นรูปธรรมที่สุด

\begin{table}[htbp]
\centering
\footnotesize\sansall
\caption{ไฟล์ที่ \en{backend} ต้องมีเพิ่มเพื่อให้ล็อกอินได้ อ้างจาก \code{spike-option-b/backend/src/}}
\label{tab:befiles}
\begin{tabular}{@{}L{4.0cm}L{9.0cm}@{}}
\toprule
\textbf{ไฟล์} & \textbf{หน้าที่ และเหตุผลที่ต้องแยกไฟล์} \\
\midrule
\code{auth/auth.router.ts} & นิยาม 3 \en{procedure} พร้อม \en{schema} ---
เป็นไฟล์เดียวที่เครื่องมืออ่านเพื่อสร้างสัญญา \\
\addlinespace[2pt]
\code{auth/auth.service.ts} & ตรรกะการล็อกอิน ไม่รู้จัก \en{tRPC} คุกกี้ หรือ \en{HTTP}
เลย จึงทดสอบได้โดยไม่ต้องยิงคำขอ \\
\addlinespace[2pt]
\code{auth/password.ts} & \en{hash} และตรวจรหัสผ่านด้วย \en{scrypt} \\
\addlinespace[2pt]
\code{auth/session.ts} & ออกและตรวจ \en{token} + ตั้งค่าคุกกี้ \\
\addlinespace[2pt]
\code{auth/auth.module.ts} & ประกอบสามไฟล์ข้างบนเข้า \en{Nest} \\
\addlinespace[2pt]
\code{trpc/context.ts} & แปลงคุกกี้ในคำขอเป็น ``ผู้ใช้คนนี้คือใคร'' \\
\addlinespace[2pt]
\code{trpc/auth.middleware.ts} & ด่านที่ปฏิเสธคำขอที่ยังไม่ล็อกอิน \\
\addlinespace[2pt]
\code{common/user.mapper.ts} & แปลงแถว \en{Prisma} เป็นรูปร่างตามสัญญา \\
\bottomrule
\end{tabular}
\end{table}

\begin{keybox}[title={ไฟล์ที่สำคัญที่สุดในตารางคือ \code{user.mapper.ts}}]
เป็นที่เดียวในระบบที่รู้จักทั้ง ``รูปร่างของ \en{Prisma}'' และ ``รูปร่างของสัญญา''
ทุก \en{procedure} ที่ส่งข้อมูลผู้ใช้ออกไปต้องผ่านที่นี่

เหตุผลไม่ใช่ความเป็นระเบียบ แต่เป็นความปลอดภัย: ตาราง \code{AccountInfo}
เก็บ \code{HashedPassword} ไว้ในแถวเดียวกัน การคืนแถว \en{Prisma}
ออกไปตรง ๆ จะทำให้รหัสผ่านที่ \en{hash} แล้วหลุดไปถึงเบราว์เซอร์
\end{keybox}

\section{บทเรียนจากโค้ดจริง --- การกันการเดาบัญชี}

หัวข้อนี้ไม่ใช่จุดที่ต้องแก้ แต่เป็นสิ่งที่ทีมควรรู้ก่อนเขียน \code{auth.service.ts} เอง
เพราะเป็นความผิดพลาดที่เขียนเองแล้วพลาดกันบ่อยมาก

ถ้าเขียนตรงไปตรงมา โค้ดจะคืนค่าทันทีเมื่อหาบัญชีไม่เจอ

\begin{lstlisting}
// the obvious version - and it leaks information
const account = await findAccount(identifier);
if (!account) throw new TRPCError({ code: 'UNAUTHORIZED' });
const ok = await verifyPassword(input.password, account.HashedPassword);
\end{lstlisting}

ปัญหาคือกรณี ``ไม่มีบัญชีนี้'' จะตอบ\textbf{เร็วกว่า}กรณี ``รหัสผิด'' อย่างเห็นได้ชัด
เพราะไม่ได้รัน \en{scrypt} ซึ่งจงใจให้ช้า คนที่ยิงอีเมลไปเรื่อย ๆ
แล้วจับเวลาจะรู้ได้ว่าอีเมลไหนมีบัญชีอยู่จริง

วิธีที่ \en{spike} ใช้คือรัน \en{scrypt} กับค่าหลอกเสมอ ให้เวลาเท่ากันทั้งสองกรณี

\begin{lstlisting}
const stored = account?.HashedPassword ?? DUMMY_HASH;
const ok = await verifyPassword(input.password, stored);
if (!account || !ok) throw new TRPCError({ code: 'UNAUTHORIZED',
                                           message: 'INVALID_CREDENTIALS' });
\end{lstlisting}

สังเกตว่าข้อความที่ตอบกลับเหมือนกันทั้งสองกรณีโดยตั้งใจ ---
ห้ามบอกผู้ใช้ว่า ``ไม่มีบัญชีนี้''

\section{สรุปสถานีที่ 4}

\begin{table}[htbp]
\centering
\footnotesize\sansall
\caption{จุดเดียวของสถานีที่ 4 --- แต่เป็นจุดที่ใหญ่ที่สุดในเอกสารเมื่อวัดเป็นปริมาณงาน}
\label{tab:station4}
\begin{tabular}{@{}cL{6.4cm}cL{3.8cm}@{}}
\toprule
\textbf{ระดับ} & \textbf{จุด} & \textbf{\en{spike} แก้แล้ว} & \textbf{ที่ยกมาใช้ได้} \\
\midrule
\pzero & \en{backend} ยังไม่มีระบบล็อกอินเลย ต้องเขียนเพิ่ม 8 ไฟล์ & ใช่ (เขียนครบและรันได้) & \code{spike-option-b/backend/src/} \\
\bottomrule
\end{tabular}
\end{table}

\clearpage

% =============================================================================
\part{สถานีที่ 5 --- ฐานข้อมูล}
% =============================================================================

นี่คือภาคที่สำคัญที่สุดของเอกสาร ทุกจุดในภาคนี้มีลักษณะร่วมกันสองอย่าง:
\textbf{\en{spike} แก้ไม่ได้} เพราะต้องเขียน \en{migration} ซึ่งกระทบทุกคน
และ\textbf{ยิ่งเลื่อนยิ่งแพง} เพราะทุกตารางที่ถูกใช้งานเพิ่มทำให้การแก้ยากขึ้น

\section{สภาพปัจจุบัน}

\code{backend/prisma/schema.prisma} มี 33 \en{model} และ \en{migration}
เพียงชุดเดียวคือ \code{20260805125233\_init}

\en{spike} ใช้\textbf{สำเนา}ของไฟล์นี้ ไม่ได้ชี้ไปที่ไฟล์จริง และตรวจแล้วว่า
เนื้อหาตรงกันทุกบรรทัด ต่างกันแค่คอมเมนต์ที่เติมไว้ตอนต้นไฟล์
ข้อค้นพบทั้งหมดในภาคนี้จึงใช้กับฐานข้อมูลจริงของทีมได้โดยตรง

\section{\pzero\ จุดที่ 14 --- คอลัมน์ที่ใช้ล็อกอินไม่มี \code{@unique}}

\begin{badbox}[title={ช่องโหว่ที่ร้ายแรงที่สุดในเอกสารนี้ · \code{schema.prisma:22, 24}}]
\begin{lstlisting}[language={}]
model AccountInfo {
  AccountKey     Int    @id @default(autoincrement())
  Email          String        // <-- no @unique
  HashedPassword String
  UserID         String        // <-- no @unique
  ...
}
\end{lstlisting}
\vspace{-0.4em}
ฐานข้อมูล\textbf{อนุญาตให้มีอีเมลซ้ำกันได้} และอนุญาตให้มี \code{UserID}
ซ้ำกันได้ ทั้งสองคอลัมน์คือคอลัมน์ที่ระบบล็อกอินใช้ค้นหาผู้ใช้
\end{badbox}

ผลที่ตามมาในโค้ดเห็นได้ทันที: \code{auth.service.ts} ต้องใช้ \code{findFirst}
ไม่ใช่ \code{findUnique} เพราะ \en{Prisma} ไม่ยอมให้ค้นแบบ \en{unique}
กับคอลัมน์ที่ไม่ได้ประกาศว่า \en{unique}

และ ``\en{first}'' แปลว่า\textbf{แถวไหนก็ได้ที่เจอก่อน}

\begin{figure}[htbp]
\centering
\begin{tikzpicture}[font=\sansall\scriptsize]

\node[blk, fill=white, draw=rule, text width=5.2cm, align=left] (t) at (0,0)
  {\textbf{ตาราง \en{AccountInfo}}\\[3pt]
   \begin{tabular}{@{}lll@{}}
   \en{Key} & \en{Email} & \en{Role}\\
   1 & \en{somchai@ku.th} & \en{Student}\\
   \textcolor{bad}{2} & \textcolor{bad}{\en{somchai@ku.th}} & \textcolor{bad}{\en{Admin}}\\
   \end{tabular}};

% เว้นระยะแนวตั้งให้พอ มิฉะนั้นกล่องสองใบจะชนกันจนอ่านลูกศรไม่ออก
\node[bestage, text width=5.0cm, align=left] (q) at (8.2,1.25)
  {\en{findFirst(\{ Email: ... \})}\\[1pt]
   \scriptsize คิวรีที่ระบบล็อกอินใช้};

\node[blk, fill=bad!10, draw=bad!50, text width=5.0cm, align=left] (r) at (8.2,-1.5)
  {\textbf{คืนแถวไหน?}\\[1pt]
   ไม่มีใครรับประกัน ลำดับขึ้นกับ\\แผนคิวรีของฐานข้อมูลในขณะนั้น};

\draw[ar] (t) -- (q);
\draw[ar] (q) -- node[lbl, right] {เจอสองแถว} (r);

% ต้องกำหนด text width ให้ node นี้ มิฉะนั้นข้อความไทยยาวจะไม่ตัดบรรทัด
% แล้วดันขอบรูปล้นออกนอกหน้ากระดาษ (overfull hbox 38pt)
\node[lbl, text=bad, align=left, anchor=north west, text width=13.8cm]
  at (-2.6,-2.9)
  {\textbf{สถานการณ์โจมตี} สมัครบัญชีด้วยอีเมลของคนอื่น
   แล้วรอให้ระบบเลือกแถวของเรา --- หรือแย่กว่านั้น
   ให้ระบบเลือกแถวของ \en{Admin} ตอนเจ้าตัวล็อกอิน};

\end{tikzpicture}
\caption{ผลของการไม่มี \code{@unique} บนคอลัมน์ที่ใช้ล็อกอิน --- ปัญหานี้ไม่ได้อยู่ที่โค้ด แต่อยู่ที่โครงสร้างตาราง}
\label{fig:unique}
\end{figure}

\begin{keybox}[title={\en{migration} ที่ต้องเขียน ข้อที่ 1}]
เติม \code{@unique} ให้ \code{Email} และ \code{UserID}

ต้องทำก่อนมีข้อมูลจริง เพราะถ้ามีแถวซ้ำอยู่แล้ว \en{migration} จะรันไม่ผ่าน
และต้องมาไล่ลบข้อมูลทีหลัง ตอนนี้ฐานข้อมูลยังว่าง --- \textbf{นี่คือช่วงเวลาที่ถูกที่สุด}
\end{keybox}

\section{\pzero\ จุดที่ 15 --- ไม่มีตารางเก็บ \en{session}}

ค้นทั้ง 33 \en{model} ไม่พบตารางชื่อ \code{Session}, \code{Token},
\code{RefreshToken} หรืออะไรที่ทำหน้าที่นี้เลยแม้แต่ตารางเดียว

นั่นทำให้เหลือทางเลือกสองทาง และทั้งสองทางมีราคา

\begin{table}[htbp]
\centering
\footnotesize\sansall
\caption{สองทางเลือกเรื่อง \en{session} และสิ่งที่ต้องแลก}
\label{tab:session}
\begin{tabular}{@{}L{3.4cm}L{4.6cm}L{4.6cm}@{}}
\toprule
& \textbf{ก. \en{JWT} ไร้สถานะ} & \textbf{ข. เพิ่มตาราง \en{Session}} \\
\midrule
\en{spike} ใช้ & ใช่ & ไม่ \\
\addlinespace[2pt]
ต้อง \en{migrate} & ไม่ต้อง & ต้อง --- แตะงานคนอื่น \\
\addlinespace[2pt]
ทำได้เมื่อไร & วันนี้ & ต้องตกลงกันก่อน \\
\addlinespace[2pt]
เพิกถอน \en{session} & \textbf{ทำไม่ได้จริง} --- ``ออกจากระบบ''
แค่ลบคุกกี้ในเบราว์เซอร์ \en{token} นั้นยังใช้ได้จนหมดอายุ &
ทำได้ และดูได้ว่าผู้ใช้ล็อกอินจากที่ไหนบ้าง \\
\addlinespace[2pt]
ต้นทุนต่อคำขอ & ไม่มีคิวรีเพิ่ม & คิวรีเพิ่ม 1 ครั้งทุกคำขอ \\
\bottomrule
\end{tabular}
\end{table}

\begin{warnbox}[title={ประโยคที่ต้องเข้าที่ประชุม}]
ถ้าเลือกทาง ก. ต้องยอมรับอย่างเป็นทางการว่า
\textbf{ถ้ามีคนขโมย \en{token} ไปได้ เราหยุดมันไม่ได้} นอกจากเปลี่ยน
\code{JWT\_SECRET} ซึ่งจะเตะผู้ใช้ทุกคนออกจากระบบพร้อมกัน

ระบบนี้ให้ยืมครุภัณฑ์ที่มีมูลค่าและมีขั้นตอนอนุมัติ การยอมรับข้อนี้
ต้องเป็นการตัดสินใจร่วม ไม่ใช่ผลพลอยได้จากการที่ไม่มีใครสังเกตว่าไม่มีตาราง
\end{warnbox}

\section{\pone\ จุดที่ 16 --- จำนวนวันที่ยืมได้ ไม่ได้ขึ้นกับเครดิตอย่างเดียว}

\code{frontend/src/constants/index.ts:44--49} เขียนตารางนี้ไว้

\begin{lstlisting}
export const CREDIT_BANDS = [
  { band: "D0", min: 80, max: 100, loanDays: 14, label: "..." },
  { band: "D1", min: 50, max:  79, loanDays:  7, label: "..." },
  { band: "D2", min: 30, max:  49, loanDays:  7, label: "..." },
  { band: "D3", min:  0, max:  29, loanDays:  5, label: "..." },
];
\end{lstlisting}

แต่ฐานข้อมูลบอกอีกอย่าง

\begin{lstlisting}[language={}]
model BorrowConstraints {
  BorrowRuleKey Int
  CreditTierKey Int
  MaxBorrowDate Int
  @@unique([BorrowRuleKey, CreditTierKey])   // <-- a PAIR, not one key
}
\end{lstlisting}

\code{MaxBorrowDate} ผูกกับ\textbf{คู่} (กฎการยืมของสิ่งของ $\times$ ระดับเครดิต)
ไม่ใช่ระดับเครดิตอย่างเดียว

\begin{badbox}[title={คำถามที่ไม่มีคำตอบเดียว}]
``ผู้ใช้คนนี้ยืมได้กี่วัน'' \textbf{ถ้าไม่บอกว่าจะยืมอะไร เป็นคำถามที่ตอบไม่ได้}

ตาราง \code{CREDIT\_BANDS} จึงผิดมาตั้งแต่บรรทัดแรกที่เขียน และผิดแบบเงียบ
เพราะข้อมูลจำลองก็ตอบตามตารางนั้นอยู่แล้ว
\end{badbox}

\en{spike} แก้ปัญหาเฉพาะหน้าโดยให้เซิร์ฟเวอร์ส่ง \code{maxBorrowDays}
ที่เป็น\textbf{ค่าสูงสุดที่เป็นไปได้} มาให้แสดงคร่าว ๆ พร้อมคอมเมนต์กำกับว่า
ตัวเลขที่ผูกพันจริงต้องมาจากตอนสร้างคำขอยืม เมื่อรู้แล้วว่าจะยืมชิ้นไหน

\begin{goodbox}[title={ทดสอบแล้วว่าตัวเลขมาจากฐานข้อมูลจริง}]
บัญชี \code{test\_lowcredit} เครดิต 20 ได้ \code{creditBand: "D3"}
และ \code{maxBorrowDays: 5} ซึ่งอ่านมาจากตาราง ไม่ใช่จากค่าคงที่ในโค้ดหน้าเว็บ
\end{goodbox}

\section{\pone\ จุดที่ 17 --- ``สังกัด'' ของผู้ใช้เป็นการเดา}

\code{AccountInfo} ไม่มี \code{FacultyKey} เส้นทางไปหาภาควิชาต้องเดินผ่านสี่ตาราง

\begin{center}
\footnotesize\sansall
\code{AccountInfo} $\rightarrow$ \code{Authority[]} $\rightarrow$
\code{ManagementGroup} $\rightarrow$ \code{FacultyInfo}
\end{center}

และ \code{Authority} มี \code{@@unique([AccountKey, ManageGroupKey])}
ซึ่งแปลว่า\textbf{คนหนึ่งคนอยู่ได้หลายกลุ่ม} เช่นอาจารย์ที่ดูแลสองภาควิชา

\en{spike} เลือกเอากลุ่มแรกที่เจอ ซึ่งเป็น\textbf{การเดา}และเขียนกำกับไว้ตรง ๆ ในโค้ด
ว่าเป็นการเดา ที่ประชุมต้องนิยามว่า ``สังกัดหลัก'' คืออะไร
หรือเปลี่ยนให้สัญญาส่งกลับเป็นรายการ

\section{\ptwo\ จุดที่ 18 --- ชื่อระดับเครดิตเป็นค่าที่ไม่บังคับ}

\code{CreditTier.CreditTierName} ประกาศเป็น \code{String?} คือมีหรือไม่มีก็ได้

เมื่อไม่มีชื่อ โค้ดต้องเดาระดับจากคะแนนแทน ซึ่งทำให้\textbf{กฎการแบ่งระดับ
ไปอยู่สองที่}: ในฐานข้อมูล และในโค้ดที่เดาแทนเมื่อฐานข้อมูลไม่บอก

ทางแก้คือทำให้ฟิลด์นี้บังคับ แล้วลบตรรกะการเดาออก

\section{\ptwo\ จุดที่ 19 --- \code{RoleInfo} เป็นตาราง ไม่ใช่ \en{enum}}

\begin{lstlisting}[language={}]
model RoleInfo {
  RoleKey  Int    @id @default(autoincrement())
  RoleName String
}
\end{lstlisting}

\en{schema} ทั้งไฟล์มี 0 \en{enum} แปลว่าชื่อบทบาทเป็นข้อความอิสระในฐานข้อมูล
โค้ดต้องมีตัวแปลงจากข้อความเป็นบทบาทตามสัญญา และต้องตัดสินใจว่า
ถ้าเจอชื่อที่ไม่รู้จักจะทำอย่างไร

\en{spike} เลือกให้\textbf{โยนข้อผิดพลาดทันที}แทนที่จะถอยไปใช้ค่าเริ่มต้น
เพราะการถอยไปใช้ค่าเริ่มต้นแปลว่าให้สิทธิ์ผิดโดยไม่มีใครรู้

\begin{lstlisting}
default:
  throw new Error(`Unknown RoleName "${roleName}" - add a mapping`);
\end{lstlisting}

\section{สรุปสถานีที่ 5}

\begin{table}[htbp]
\centering
\footnotesize\sansall
\caption{หกจุดของสถานีที่ 5 --- ไม่มีข้อไหนที่ \en{spike} แก้ได้ เพราะทุกข้อต้องตกลงกันก่อน}
\label{tab:station5}
\begin{tabular}{@{}cL{6.8cm}L{4.4cm}@{}}
\toprule
\textbf{ระดับ} & \textbf{จุด} & \textbf{ต้องทำอะไร} \\
\midrule
\pzero & \code{Email}/\code{UserID} ไม่มี \code{@unique} & \en{migration} \\
\pzero & ไม่มีตารางเก็บ \en{session} & ตัดสินใจ แล้วอาจ \en{migrate} \\
\pone & \code{maxBorrowDays} ขึ้นกับคู่ ไม่ใช่เครดิตเดี่ยว & แก้สัญญา + ลบ \code{CREDIT\_BANDS} \\
\pone & ``สังกัดหลัก'' ยังไม่ถูกนิยาม & ตัดสินใจ \\
\ptwo & \code{CreditTierName} ไม่บังคับ & \en{migration} เล็ก \\
\ptwo & \code{RoleInfo} เป็นตาราง & เขียนตัวแปลง + ตกลงรายชื่อบทบาท \\
\bottomrule
\end{tabular}
\end{table}

\clearpage

% =============================================================================
\part{ต้องแก้อะไรบ้าง และเรียงลำดับอย่างไร}
% =============================================================================

\section{ตารางรวมทั้ง 19 จุด}

\begin{table}[htbp]
\centering
\footnotesize\sansall
\caption{ทุกจุดที่พบ เรียงตามความรุนแรง --- คอลัมน์สุดท้ายคือคำถามที่สำคัญที่สุด}
\label{tab:all}
\begin{tabular}{@{}cL{5.5cm}L{2.5cm}L{3.4cm}@{}}
\toprule
\textbf{ระดับ} & \textbf{จุด} & \textbf{สถานี} & \textbf{\en{spike} แก้ให้แล้วหรือยัง} \\
\midrule
\pzero & \code{Email}/\code{UserID} ไม่มี \code{@unique} & ฐานข้อมูล & \textbf{ไม่ --- ต้องตกลงกัน} \\
\pzero & ไม่มีตารางเก็บ \en{session} & ฐานข้อมูล & \textbf{ไม่ --- ต้องตกลงกัน} \\
\pzero & บทบาทเก็บใน \en{localStorage} & หน้าจอ & ใช่ \\
\pzero & รหัสผ่านทดสอบอยู่ในรีโป & หน้าจอ & ใช่ \\
\pzero & \en{backend} ไม่มีระบบล็อกอิน (8 ไฟล์) & \en{server} & ใช่ ยกไปใช้ได้ \\
\midrule
\pone & \code{zod} 3 ต้องยกเป็น 4 & สัญญา & ใช่ \\
\pone & \code{User} เขียนมือ ผิด 3 ฟิลด์ & สัญญา & ใช่ \\
\pone & หน้าล็อกอินเดาบทบาทเอง & หน้าจอ & ใช่ \\
\pone & \code{main.ts} ไม่มี \en{CORS}/\en{cookie-parser} & เครือข่าย & ใช่ ยกไปใช้ได้ \\
\pone & \code{maxBorrowDays} ขึ้นกับคู่ & ฐานข้อมูล & บรรเทาชั่วคราว \\
\pone & ``สังกัดหลัก'' ยังไม่ถูกนิยาม & ฐานข้อมูล & \textbf{ไม่ --- ต้องตกลงกัน} \\
\midrule
\ptwo & รหัสข้อผิดพลาดคนละชุด & สัญญา & ใช่ \\
\ptwo & \en{CI} มองไม่เห็นผู้บริโภครายที่สอง & สัญญา & \textbf{ไม่} \\
\ptwo & ฟอร์มต้องเป็น \en{async} & หน้าจอ & ใช่ (2 ไฟล์) \\
\ptwo & แถบข้างแสดงชื่อตายตัว & หน้าจอ & \textbf{ไม่} \\
\ptwo & แยกกรณีต่อไม่ติดไม่ได้ & เครือข่าย & ใช่ \\
\ptwo & \en{proxy} \code{/api} ชี้ \code{:3000} & เครือข่าย & \textbf{ไม่ --- ต้องตัดสินใจ} \\
\ptwo & \code{CreditTierName} ไม่บังคับ & ฐานข้อมูล & \textbf{ไม่} \\
\ptwo & \code{RoleInfo} เป็นตาราง & ฐานข้อมูล & บรรเทาแล้ว (มีตัวแปลง) \\
\bottomrule
\end{tabular}
\end{table}

\begin{keybox}[title={อ่านตารางนี้แล้วจะเห็นรูปแบบหนึ่ง}]
ทุกจุดที่ \en{spike} \textbf{แก้ให้แล้ว} คือจุดที่แก้ได้ด้วยคนคนเดียวในไฟล์ของตัวเอง

ทุกจุดที่ \en{spike} \textbf{แก้ไม่ได้} คือจุดที่ต้องแตะของกลาง ---
\en{schema} ฐานข้อมูล หรือข้อตกลงของทีม

นั่นแปลว่า\textbf{งานที่เหลือส่วนใหญ่ไม่ใช่งานเขียนโค้ด แต่เป็นงานตัดสินใจ}
และการตัดสินใจเลื่อนได้เรื่อย ๆ จนกว่าจะมีคนกำหนดวัน
\end{keybox}

\section{ลำดับที่ต้องลงมือ}

ลำดับนี้ไม่ได้เรียงตามความสำคัญอย่างเดียว แต่เรียงตาม\textbf{การพึ่งพากัน}
งานบางอย่างทำก่อนแล้วถูกกว่ามาก

\begin{figure}[htbp]
\centering
\begin{tikzpicture}[font=\sansall\scriptsize, node distance=6mm]

\node[blk, fill=bad!10, draw=bad!50, text width=3.5cm] (d1)
  {\textbf{ขั้นที่ 1 · ตัดสินใจ}\\ตาราง \en{Session} เอาหรือไม่\\นิยาม ``สังกัดหลัก''};
\node[blk, fill=bad!10, draw=bad!50, text width=3.5cm, right=8mm of d1] (d2)
  {\textbf{ขั้นที่ 2 · \en{migration}}\\เติม \code{@unique}\\(ทำตอนฐานข้อมูลยังว่าง)};
\node[stage, text width=3.5cm, right=8mm of d2] (d3)
  {\textbf{ขั้นที่ 3 · ยก \code{zod}}\\3 $\rightarrow$ 4 ใน \code{frontend/}\\ตอนนี้ยังฟรี};

\node[bestage, text width=3.5cm, below=10mm of d1] (d4)
  {\textbf{ขั้นที่ 4 · \en{backend}}\\ยก 8 ไฟล์จาก \en{spike}\\+ ตั้ง \en{CORS}};
\node[festage, text width=3.5cm, right=8mm of d4] (d5)
  {\textbf{ขั้นที่ 5 · \en{frontend}}\\ยก 14 ไฟล์ที่แก้แล้ว\\ลบ \code{mock-auth.ts}};
\node[opstage, text width=3.5cm, right=8mm of d5] (d6)
  {\textbf{ขั้นที่ 6 · \en{CI}}\\แก้ให้วนทุกโฟลเดอร์\\ที่กินสัญญา};

\draw[ar] (d1) -- (d2);
\draw[ar] (d2) -- (d3);
\draw[ar] (d3.south) -- ++(0,-0.5) -| (d4.north);
\draw[ar] (d4) -- (d5);
\draw[ar] (d5) -- (d6);

\end{tikzpicture}
\caption{ลำดับการลงมือที่เรียงตามการพึ่งพากัน ไม่ใช่ตามความสำคัญเพียงอย่างเดียว}
\label{fig:order}
\end{figure}

สองประโยคที่อธิบายรูปที่~\ref{fig:order} ได้ครบ

\begin{itemize}
  \item ขั้นที่ 1--2 คืองานที่\textbf{ต้องทำก่อนใครเขียนโค้ดเพิ่ม}
        เพราะเป็นการแก้ของกลาง ยิ่งมีคนสร้างงานทับข้างบนมากเท่าไร การรื้อยิ่งแพง
  \item ขั้นที่ 3 ทำตอนนี้\textbf{ฟรี} ทำทีหลังจ่ายตามจำนวนฟอร์มที่เพิ่มเข้ามาระหว่างนั้น
\end{itemize}

\begin{warnbox}[title={ขั้นที่ 2 มีหน้าต่างเวลาที่กำลังจะปิด}]
การเติม \code{@unique} ตอนฐานข้อมูลยังว่างคือการแก้ไฟล์บรรทัดเดียวแล้วรัน \en{migration}

การเติม \code{@unique} ตอนมีข้อมูลจริงแล้วคือการต้องไล่หาแถวซ้ำ ตัดสินใจว่าจะลบแถวไหน
ติดต่อผู้ใช้ที่ได้รับผลกระทบ แล้วค่อยรัน \en{migration}

\textbf{ตอนนี้ฐานข้อมูลยังว่าง} นี่คือช่วงเวลาที่ถูกที่สุดและกำลังหมดลง
\end{warnbox}

\section{สิ่งที่ยังไม่รู้}

เอกสารนี้พยายามไม่เดา สิ่งที่ยังไม่มีคำตอบขอเขียนไว้ตรง ๆ

\begin{itemize}
  \item \textbf{ต่อได้แค่ \en{auth}} --- \en{backend} มี 3 \en{procedure}
        จากที่ระบบต้องการทั้งหมดราว 18 ยังไม่รู้ว่า \en{procedure} ที่ซับซ้อนกว่านี้
        (เช่นการสร้างคำขอยืมที่มีหลายรายการ) จะเจอข้อจำกัดอะไรเพิ่ม
  \item \textbf{ยังไม่ได้ทดสอบตอน \en{session} หมดอายุกลางคัน} --- ตั้ง \en{TTL} ไว้ 8 ชั่วโมง
        แต่ยังไม่ได้ลองว่าหน้าจอจะทำตัวอย่างไรเมื่อหมดอายุระหว่างใช้งาน
  \item \textbf{ยังไม่ได้ทดสอบบนมือถือและยังไม่ได้ตรวจ \en{accessibility}}
  \item \textbf{ยังไม่ได้ลองรัน \en{CI} บน \en{GitHub} จริง} ---
        สคริปต์ทุกตัวทดสอบแล้วในเครื่อง แต่ตัว \en{workflow} เองยังไม่เคยผ่าน \en{runner}
        ยังไม่รู้ว่าไบนารีของเครื่องมือสร้างสัญญามีรุ่นที่ตรงกับ \en{ubuntu-latest} หรือไม่
\end{itemize}

\clearpage

% =============================================================================
\appendix
\part{ภาคผนวก}
% =============================================================================

\section*{ภาคผนวก ก --- ทุกไฟล์ที่ \en{spike} แก้}
\addcontentsline{toc}{section}{ภาคผนวก ก --- ทุกไฟล์ที่ \en{spike} แก้}

ตารางนี้คือคำตอบครบถ้วนของคำถาม ``\en{spike} แก้อะไรจากโค้ดจริงบ้าง''
นับจาก \code{diff} ระหว่าง \code{frontend/} กับ \code{spike-option-b/frontend-real/}

\begin{table}[htbp]
\centering
\footnotesize\sansall
\caption{ไฟล์ที่แก้ 13 ไฟล์ --- ตัวเลขคือบรรทัดที่เพิ่ม/ลบ}
\label{tab:diffmod}
\begin{tabular}{@{}L{4.6cm}ccL{5.6cm}@{}}
\toprule
\textbf{ไฟล์} & \textbf{+} & \textbf{--} & \textbf{แก้อะไร} \\
\midrule
\code{package.json} & 5 & 3 & \code{zod} 3$\rightarrow$4.4.3, \code{resolvers} 3$\rightarrow$5.2.2, เพิ่ม \code{@trpc/client} \\
\code{vite.config.ts} & 7 & 8 & พอร์ต 5173$\rightarrow$5275, ตัด \en{proxy} \\
\code{src/types/domain.ts} & 18 & 12 & \code{User}/\code{Role}/\code{CreditBand} ดึงจากสัญญา \\
\code{src/types/env.d.ts} & 2 & 0 & เพิ่ม \code{VITE\_TRPC\_URL} \\
\code{src/lib/error-messages.ts} & 30 & 1 & รหัสตรงกับ \en{backend} + แยกกรณีต่อไม่ติด \\
\code{src/app/app.tsx} & 28 & 18 & \code{hydrate()} $\rightarrow$ \code{<SessionBootstrap>} \\
\code{src/features/auth/auth.store.ts} & 19 & 38 & ตัด \code{loginAs} และ \code{hydrate} ทิ้ง \\
\code{src/features/auth/use-me.ts} & 15 & 6 & เรียก \code{trpc.auth.me} \\
\code{src/features/auth/login-page.tsx} & 34 & 24 & เรียกเซิร์ฟเวอร์จริง บทบาทมาจากคำตอบ \\
\code{...login-method-ku.tsx} & 24 & 6 & \en{async} + ปุ่มมีสถานะ + ที่แสดงข้อผิดพลาด \\
\code{...login-method-local.tsx} & 14 & 8 & \en{async} + ปุ่มมีสถานะ \\
\code{src/components/layout/sidebar.tsx} & 8 & 3 & \en{logout} เป็นการเรียกเซิร์ฟเวอร์ \\
\code{src/i18n/locales/\{th,en\}.json} & 4 & 2 & เพิ่มคำว่า ``กำลังเข้าสู่ระบบ'' \\
\bottomrule
\end{tabular}
\end{table}

\begin{table}[htbp]
\centering
\footnotesize\sansall
\caption{ไฟล์ใหม่ 4 ไฟล์ และไฟล์ที่ลบ 1 ไฟล์}
\label{tab:diffnew}
\begin{tabular}{@{}L{5.4cm}cL{6.4cm}@{}}
\toprule
\textbf{ไฟล์} & \textbf{บรรทัด} & \textbf{คืออะไร} \\
\midrule
\code{src/lib/trpc.ts} & 66 & ตัวเรียก \en{tRPC} + ตัวแยกประเภทข้อผิดพลาด \\
\code{src/generated/server.ts} & 56 & \textbf{สัญญา} --- สร้างโดยเครื่องมือ ห้ามแก้มือ \\
\code{src/features/auth/use-login.ts} & 28 & \en{mutation} สำหรับล็อกอิน \\
\code{src/features/auth/use-logout.ts} & 29 & \en{mutation} สำหรับออกจากระบบ \\
\midrule
\code{src/features/auth/mock-auth.ts} & \textbf{--70} & \textbf{ลบทั้งไฟล์} \\
\bottomrule
\end{tabular}
\end{table}

\begin{keybox}[title={ตัวเลขรวมที่ควรจำ}]
แก้ 13 ไฟล์ (เพิ่ม 208 บรรทัด ลบ 129) เพิ่มใหม่ 4 ไฟล์ (179 บรรทัด)
ลบทิ้ง 1 ไฟล์ (70 บรรทัด)

เทียบกับขนาดของ \code{frontend/src/} ทั้งหมดแล้วนี่คือ\textbf{ส่วนน้อยมาก}
--- และไม่มีไฟล์ในกลุ่ม \code{components/ui/}, \code{features/admin/}
หรือ \code{app/router.tsx} ถูกแตะเลยแม้แต่ไฟล์เดียว
\end{keybox}

\section*{ภาคผนวก ข --- วิธีรัน \en{spike} เพื่อดูด้วยตาเอง}
\addcontentsline{toc}{section}{ภาคผนวก ข --- วิธีรัน \en{spike} เพื่อดูด้วยตาเอง}

\begin{lstlisting}[language=bash]
# 1. database
docker compose -f spike-option-b/docker-compose.yml up -d

# 2. backend on :3002
cd spike-option-b/backend && npm ci && cp .env.example .env
npx prisma generate && npm run dev

# 3. the team's real frontend, wired, on :5275
cd spike-option-b/frontend-real && npm ci && npm run dev
\end{lstlisting}

เปิด \code{http://localhost:5275} แล้วใช้บัญชีในตารางที่~\ref{tab:accounts}
รหัสผ่านทั้งหมดถูก \en{hash} ด้วย \en{scrypt} เก็บในฐานข้อมูลจริง

\begin{table}[htbp]
\centering
\footnotesize\sansall
\caption{บัญชีทดสอบที่ \code{seed.ts} สร้างไว้}
\label{tab:accounts}
\begin{tabular}{@{}L{2.6cm}L{3.8cm}L{2.8cm}L{3.4cm}@{}}
\toprule
\textbf{ช่องที่ใช้} & \textbf{กรอก} & \textbf{รหัสผ่าน} & \textbf{บทบาทที่ได้} \\
\midrule
อีเมล \en{KU} & \code{natthawut.s@ku.th} & \code{borrower1234} & \en{borrower} \\
อีเมล \en{KU} & \code{orawan.p@ku.th} & \code{supervisor1234} & \textbf{\en{supervisor}} \\
อีเมล \en{KU} & \code{lowcredit.s@ku.th} & \code{lowcredit1234} & \en{borrower} เครดิต 20 \\
บัญชีภายใน & \code{test\_staff} & \code{staff1234} & \en{staff} \\
บัญชีภายใน & \code{test\_admin} & \code{admin1234} & \en{admin} \\
\bottomrule
\end{tabular}
\end{table}

สิ่งที่ควรลองกดมี 4 อย่างที่เห็นผลชัดที่สุด

\begin{enumerate}
  \item ล็อกอินด้วย \code{orawan.p@ku.th} --- ต้องไปโผล่หน้าอนุมัติของหัวหน้า
        ไม่ใช่หน้าผู้ยืม (จุดที่ 1)
  \item ล็อกอินแล้วกด \en{F5} --- ต้องยังอยู่ในระบบ (สถานีที่ 3)
  \item เปิดเครื่องมือนักพัฒนา พิมพ์ \code{document.cookie} --- ต้องไม่เห็น
        \code{ulms\_session} (จุดที่ 2)
  \item ปิด \en{backend} แล้วลองล็อกอิน --- ต้องขึ้นข้อความว่าเชื่อมต่อไม่ได้
        ไม่ใช่ว่ารหัสผ่านผิด (จุดที่ 12)
\end{enumerate}

\section*{ภาคผนวก ค --- ผลการทดสอบ 28 ข้อ}
\addcontentsline{toc}{section}{ภาคผนวก ค --- ผลการทดสอบ 28 ข้อ}

\subsection*{ชั้นข้อมูล 13 ข้อ}
\addcontentsline{toc}{subsection}{ชั้นข้อมูล 13 ข้อ}

รันด้วย \code{node spike-option-b/frontend-real/e2e-check.mjs}
ใช้ \code{@trpc/client} ตัวเดียวกับที่แอปใช้ จึงครอบคลุมรูปแบบการส่งข้อมูลจริง
และวิธีที่ข้อผิดพลาดจากเซิร์ฟเวอร์กลายเป็น \code{TRPCClientError} ฝั่งหน้าเว็บ

\begin{lstlisting}[language=bash]
PASS  auth.me without a session -> NOT_AUTHENTICATED 401
PASS  auth.login local test_staff -> role staff
PASS    contract field id is a number, not a string
PASS    contract field maxBorrowDays present
PASS  auth.me after login -> same user
PASS  auth.login wrong password -> INVALID_CREDENTIALS 401
PASS  auth.login non-KU email -> NOT_KU_EMAIL 400
PASS  auth.login identifier > 200 chars -> BAD_REQUEST
PASS  auth.login KU email of an academic -> role supervisor
PASS  low-credit student -> D3 from the DB
PASS    maxBorrowDays for D3 comes from the DB, not CREDIT_BANDS
PASS  auth.logout -> ok
PASS  auth.me after logout -> NOT_AUTHENTICATED
\end{lstlisting}

\subsection*{ชั้นหน้าจอ 15 ข้อ}
\addcontentsline{toc}{subsection}{ชั้นหน้าจอ 15 ข้อ}

รันด้วย \code{node spike-option-b/browser-check/ui-check.mjs}
เปิด \en{Firefox} จริง กรอกฟอร์มจริง

\begin{lstlisting}[language=bash]
PASS  visiting / with no session lands on /login
PASS  wrong password shows the mapped Thai message
PASS    and stays on /login
PASS  staff login lands on the staff dashboard
PASS  browser stored the session cookie
PASS    and it is httpOnly (JS cannot read it)
PASS    document.cookie really cannot see it
PASS  reload keeps the session (auth.me bootstrap)
PASS  staff visiting /admin is bounced back to /staff
PASS  logout returns to /login
PASS    session cookie is gone
PASS    and /staff is no longer reachable
PASS  academic signing in with a KU email lands on SUPERVISOR home
PASS  no console errors during the whole run
\end{lstlisting}

\begin{keybox}[title={สามข้อในรายการบนที่ทดสอบด้วยวิธีอื่นไม่ได้เลย}]
\textbf{คุกกี้เป็น \code{httpOnly} จริง} --- ต้องมี \en{JavaScript}
ที่รันอยู่ในหน้าเว็บถึงจะลองอ่านได้ สคริปต์ยิง \en{HTTP} เฉย ๆ ตรวจข้อนี้ไม่ได้

\textbf{รีเฟรชแล้วไม่หลุด} --- พิสูจน์ว่ากลไกกู้ \en{session} ทำงานถูก

\textbf{\en{RBAC} ยึดจริง} --- พิมพ์ \en{URL} ตรง ๆ แล้วถูกเด้งกลับ
\end{keybox}

\section*{ภาคผนวก ง --- ที่มาของข้อเท็จจริงทุกข้อ}
\addcontentsline{toc}{section}{ภาคผนวก ง --- ที่มาของข้อเท็จจริงทุกข้อ}

ถ้าไฟล์เหล่านี้เปลี่ยน เนื้อความในเอกสารต้องตามไปแก้

\begin{table}[htbp]
\centering
\footnotesize\sansall
\caption{ข้อเท็จจริงในเอกสารนี้และแหล่งที่มา อ้างถึงบรรทัด ณ วันที่ 10 ส.ค. 2569}
\label{tab:sources}
\begin{tabular}{@{}L{6.2cm}L{6.8cm}@{}}
\toprule
\textbf{ข้อเท็จจริง} & \textbf{ไฟล์:บรรทัด} \\
\midrule
หน้าล็อกอินกำหนดบทบาทตายตัว & \code{frontend/src/features/auth/login-page.tsx:41} \\
บทบาทเก็บใน \en{localStorage} & \code{frontend/src/features/auth/auth.store.ts:36, 46} \\
ตารางรหัสผ่านทดสอบ & \code{frontend/src/features/auth/mock-auth.ts:62} \\
รหัสผ่านทดสอบใน \en{README} & \code{frontend/README.md:14--18} \\
\code{User.id} เป็น \code{string} & \code{frontend/src/types/domain.ts:10} \\
\code{departmentId} ไม่ยอมรับค่าว่าง & \code{frontend/src/types/domain.ts:15} \\
\code{zod} รุ่น 3 & \code{frontend/package.json:40} \\
ตาราง \code{CREDIT\_BANDS} & \code{frontend/src/constants/index.ts:44--49} \\
รหัส \code{INVALID\_DOMAIN} & \code{frontend/src/lib/error-messages.ts:11} \\
\en{proxy} ชี้ \code{:3000} & \code{frontend/vite.config.ts} \\
\en{backend} ไม่มี \en{CORS} & \code{backend/src/main.ts} ทั้งไฟล์ 7 บรรทัด \\
\code{Email} ไม่มี \code{@unique} & \code{backend/prisma/schema.prisma:22} \\
\code{UserID} ไม่มี \code{@unique} & \code{backend/prisma/schema.prisma:24} \\
ไม่มีตาราง \en{session} & \code{backend/prisma/schema.prisma} ทั้งไฟล์ 33 \en{model} \\
\code{MaxBorrowDate} ผูกกับคู่ & \code{backend/prisma/schema.prisma:232} \\
\code{CreditTierName} ไม่บังคับ & \code{backend/prisma/schema.prisma} \en{model} \code{CreditTier} \\
\code{RoleInfo} เป็นตาราง & \code{backend/prisma/schema.prisma} \en{model} \code{RoleInfo} \\
สถานะหน้าจอ 8/26 & \code{frontend/src/features/**/*-page.tsx} \\
\en{diff} ทุกไฟล์ & \code{spike-option-b/frontend-real/} \\
ผลทดสอบชั้นข้อมูล & \code{spike-option-b/frontend-real/e2e-check.mjs} \\
ผลทดสอบชั้นหน้าจอ & \code{spike-option-b/browser-check/ui-check.mjs} \\
ภาพหน้าจอ & \code{spike-option-b/browser-check/shots/} \\
\bottomrule
\end{tabular}
\end{table}

\vfill
\begin{center}
\footnotesize\sansall\color{muted}
รูปที่วาดด้วย \en{TikZ} เป็นภาพอธิบาย · รูปที่เป็นภาพถ่ายหน้าจอเป็นผลการทดสอบจริง\\
เอกสารนี้สร้างจาก \code{docs/spike-gap.tex} ด้วย \code{latexmk spike-gap.tex}
\end{center}

\end{document}
